% ======================================================================
\chapter{Основы статистического контроля процессов}
% ======================================================================

\section*{Обзор главы}
В этой главе мы сосредоточимся на основных понятиях, терминологии и некоторых базовых статистических теориях одномерных контрольных карт. Это включает такие темы, как контрольные границы, их расположение, объем выборки, частоту выборки и то, как формировать рациональные подгруппы. Мы начнем с классических контрольных карт Шухарта. Также будут представлены некоторые более продвинутые, но практичные темы, такие как эффективность и статистическое проектирование контрольной карты.

\section{Основные понятия}
Чтобы описать и понять лежащую в основе теорию, мы начнем с обсуждения некоторых важных терминов.

\subsection{Типы изменчивости}
В любом процессе всегда будет присутствовать определенное количество естественной изменчивости, независимо от того, насколько хорошо процесс был спроектирован, разработан или поддерживается. Например, не все автомобильные двери, окрашенные одной и той же машиной при одинаковых настройках, будут идентичными. Не все пятифунтовые пакеты картофеля, упакованные одним и тем же заводом, будут содержать ровно пять фунтов картофеля, как бы тщательно они ни были упакованы. Такая изменчивость присуща процессу и не может быть (или не должна быть) полностью устранена, возможно, не без огромных затрат при минимальной отдаче. И наоборот, в процессе могут присутствовать и другие виды изменчивости, которые являются разрушительными и могут привести к ухудшению качества. Цель статистического контроля процессов (SPC) — выявить источник(и) этой изменчивости, изучить первопричину(ы) и устранить или минимизировать их, чтобы качество продукта было улучшено и поддерживалось на удовлетворительном уровне.

Доктор Уолтер А. Шухарт признал разницу между этими двумя типами изменчивости и первым ввел терминологию общих причин и особых причин (см. Shewhart, 1926, 1931).

Термин \textbf{общие причины} (или случайные причины) относится к тем источникам изменчивости, которые присущи или естественно присутствуют в процессе. Вариация, обусловленная такими причинами, неконтролируема или постоянно присутствует и также называется «фоновым шумом» внутри процесса, который может быть вызван кумулятивным эффектом многих мелких и необнаружимых, но неизбежных причин.

\textbf{Особые причины} (или назначаемые причины) — это те источники изменчивости, которые не являются частью общих причин или естественной изменчивости процесса и могут напрямую влиять на качество процесса. В производственном процессе особые причины обычно могут возникать из-за дефектного сырья, неправильно настроенных или контролируемых машин или ошибок оператора.

Общая вариация, присутствующая в процессе, представляет собой сумму вариаций от этих двух источников, то есть вариации из-за общих причин и вариации из-за особых причин.

Процесс считается находящимся в состоянии статистической управляемости или \textbf{в статистически управляемом состоянии} (in-control, IC), если он функционирует в соответствии с тем, что является целевым или ожидаемым при наличии общих причин. Таким образом, когда процесс находится в состоянии IC, вокруг некоторого целевого уровня происходят только случайные отклонения или флуктуации. Мы пишем IC для обозначения состояния (статуса) IC процесса. Когда процесс не является IC, говорят, что он \textbf{вышел из-под контроля} (out-of-control, OOC). Такое остояние процесса обозначается как OOC. Это когда есть основания полагать, что присутствие некоторых особых причин влияет на процесс, и поэтому их необходимо выявить и, если возможно, устранить, чтобы процесс вернулся в состояние IC.

Подчеркнем, цель SPC — минимизировать и/или устранить назначаемые причины изменчивости, чтобы процесс работал только под действием общих причин и при этом соответствовал или превышал любые заданные спецификации. Для достижения этой цели обычно необходимы три вещи: (i) \textbf{идентификация}: обнаружение присутствия любых особых причин вариации по «сигналу» с карты; (ii) \textbf{отслеживание}: поиск причины и возможность проследить источник сигнала; и (iii) \textbf{действие}: способность предпринять любые корректирующие действия, необходимые для возвращения процесса в состояние IC. Контрольная карта, основанная на обоснованных статистических принципах, является графическим инструментом с визуально привлекательным интерфейсом и обеспечивает объективный способ понимания изменчивости процесса, тем самым поддерживая и улучшая качество. Таким образом, контрольная карта играет важную роль в SPC.

\subsection{Контрольная карта}
Контрольная карта — это графическое отображение, состоящее из графика последовательных значений подходящей сводной статистики. Эта статистика называется \textbf{статистикой для нанесения на карту} (charting statistic) или статистикой для построения графика, которая является оценкой интересующего параметра качества и рассчитывается по каждой из последовательных выборок измерений (данных), случайно взятых по характеристике качества по вертикальной оси, в зависимости от номера выборки (или времени) по горизонтальной оси, вместе с тремя линиями, нанесенными на дисплей, которые называются контрольными границами. Обычно контрольные границы состоят из центральной линии (CL) и двух горизонтальных линий, по одной с каждой стороны от CL. Линия над CL называется верхней контрольной границей (UCL), а линия под CL называется нижней контрольной границей (LCL). Эти три линии помогают пользователю решить, остается ли процесс в состоянии IC или переходит в состояние OOC.

Значения статистики для нанесения на карту отображаются на контрольной карте для мониторинга процесса. Чтобы найти подходящую статистику для нанесения на карту, необходимо понять проблему управления процессом, определить параметр процесса, подлежащий мониторингу, и выбрать подходящую статистику, которая часто (основана на) является оценкой этого параметра, обладающей желательными статистическими свойствами, такими как несмещенность, минимальная дисперсия и т.д., как обсуждалось в Главе 1. Обычно, когда статистика для нанесения на карту попадает на или за пределы любой из контрольных границ, это сигнализирует о возможном наличии некоторых назначаемых причин, и в этот момент процесс может быть объявлен OOC. Это событие называется \textbf{сигнальным событием} или \textbf{тревогой}. \textbf{Ложная тревога} — это когда процесс объявляется OOC, хотя на самом деле он IC. В приложениях SPC вероятность ложной тревоги называется \textbf{частотой ложных тревог} и сокращается как FAR. Как только наблюдается сигнал, обычно начинается расследование для поиска источника сигнала и, при необходимости, предпринимаются корректирующие действия. Напротив, когда статистика для нанесения на карту отображается между верхней и нижней границами процесса и не показывает систематической картины во времени, говорят, что процесс находится в состоянии IC. Другими словами, на контрольной карте не наблюдается сигнала. Соответствующее событие называется \textbf{несигнальным событием}. Типичная контрольная карта показана на рисунке \ref{fig-2.1}.

\begin{figure}[!ht]
	\centering
	\includegraphics[width=0.7\linewidth]{"./images/fig-2.1.png"}
	\caption[]{Типичная контрольная карта}
	\label{fig-2.1}
\end{figure}

CL проводится либо на ожидаемом значении статистики для нанесения на карту, когда процесс находится в состоянии IC, либо на ее оценке для интересующего параметра. UCL и LCL размещаются на некотором расстоянии выше и ниже CL таким образом, чтобы, если процесс действительно IC, вероятность того, что точка попадет на или за пределы контрольных границ, была очень мала. Степень малости зависит от конкретного контекста и серьезности последствий. Это означает, что по замыслу почти все, за очень небольшим исключением, нанесенные точки, то есть значения статистики для нанесения на карту из различных выборок, взятых из процесса, как ожидается, будут попадать между контрольными границами, когда процесс действительно IC. Когда все точки отображаются между контрольными границами, это обычно указывает на отсутствие каких-либо особых причин вариации и корректирующие действия, как правило, не требуются. Это когда процесс объявляется находящимся в состоянии статистической управляемости или просто IC. Однако обратите внимание, что даже когда все точки отображаются между контрольными границами, проблемы все еще могут быть, например, когда все точки находятся между контрольными границами, но график показывает некоторую неслучайную картину, или когда точки находятся ниже или выше целевого значения (того, что может ожидаться), может возникнуть опасение, что могут присутствовать особые (назначаемые) причины и процесс может быть OOC.

Например, на рисунке \ref{fig-2.2} показана контрольная карта Шухарта, на которой все точки попадают между UCL и LCL, а картина графика выглядит случайной. Таким образом, здесь не наблюдается сигнальных событий, и нет свидетельств каких-либо аномальных паттернов среди точек. Следовательно, можно предположить, что процесс находится в состоянии IC относительно отслеживаемого параметра, и мониторинг процесса продолжается без перерыва. Однако на рисунке \ref{fig-2.3}, хотя все точки отображаются между UCL и LCL, наблюдается явная циклическая картина. Следовательно, есть реальное опасение, что могут присутствовать некоторые назначаемые причины (т.е. процесс может быть OOC), поэтому может потребоваться поиск таких назначаемых причин. Циклические паттерны могут возникать из-за изменений окружающей среды, таких как изменения температуры. Это также может быть результатом усталости оператора, колебаний напряжения или давления, или какой-либо другой переменной в оборудовании. Свойства сырья также могут играть роль. На рисунке \ref{fig-2.4} показан другой пример, где все точки отображаются между UCL и LCL; однако точки, кажется, группируются вокруг CL. Это называется \textbf{стратификацией}, и есть опасение, что процесс может быть OOC из-за заметного отсутствия изменчивости. Стратификация может быть результатом неправильного расчета контрольных границ или неправильного разделения на подгруппы; рациональное разделение на подгруппы обсуждается далее в книге. Это также может быть результатом не пересчета контрольных границ после улучшения процесса. Рисунки \ref{fig-2.3} и \ref{fig-2.4} являются двумя примерами, где все точки отображаются между двумя контрольными границами; однако ясно, что есть проблема. Мы обсудим вопрос паттернов более подробно в следующем разделе. На рисунках \ref{fig-2.3} и \ref{fig-2.4} были проиллюстрированы две контрольные карты, где все точки отображаются между контрольными границами; однако на рисунке \ref{fig-2.5} показана контрольная карта, где одна точка отображается за пределами UCL или LCL (таким образом, имеет место сигнальное событие), указывая на то, что процесс может быть OOC.

\begin{figure}[!ht]
	\centering
	\includegraphics[width=0.7\linewidth]{"./images/fig-2.2.png"}
	\caption[]{Контрольная карта Шухарта со случайным паттерном}
	\label{fig-2.2}
\end{figure}

\begin{figure}[!ht]
	\centering
	\includegraphics[width=0.7\linewidth]{"./images/fig-2.3.png"}
	\caption[]{Контрольная карта Шухарта с циклическим паттерном}
	\label{fig-2.3}
\end{figure}

\begin{figure}[!ht]
	\centering
	\includegraphics[width=0.7\linewidth]{"./images/fig-2.4.png"}
	\caption[]{Контрольная карта Шухарта со стратифицированным паттерном}
	\label{fig-2.4}
\end{figure}

\begin{figure}[!ht]
	\centering
	\includegraphics[width=0.7\linewidth]{"./images/fig-2.5.png"}
	\caption[]{Контрольная карта Шухарта с точкой, выходящей за верхнюю контрольную границу}
	\label{fig-2.5}
\end{figure}

Обратите внимание, что принято соединять точки на контрольной карте отрезками прямых для облегчения визуализации.

Контрольные карты играют решающую роль в различных областях и отраслях по всему миру, особенно после Второй мировой войны. Монтгомери (2009) перечисляет пять ключевых причин их популярности:

\begin{itemize}
	\item Контрольные карты — это проверенная техника повышения производительности.
	\item Контрольные карты эффективны в предотвращении дефектов.
	\item Контрольные карты предотвращают ненужную настройку процесса.
	\item Контрольные карты предоставляют диагностическую информацию.
	\item Контрольные карты предоставляют информацию о возможностях процесса.
\end{itemize}

Он продолжает: «Контрольные карты являются одними из самых важных инструментов управления; они так же важны, как контроль затрат и контроль материалов. Современные компьютерные технологии сделали простым внедрение контрольных карт в любой тип процесса, поскольку сбор и анализ данных могут выполняться на микрокомпьютере или терминале локальной сети в реальном времени, онлайн на рабочем центре».

Мы начнем с того, как построить одномерную контрольную карту Шухарта. Пусть \(W\) — выборочная статистика, которая измеряет некоторую характеристику качества, представляющую интерес, и предположим, что среднее значение и дисперсия \(W\) в состоянии IC равны \(\mu_w\) и \(\sigma_w^2\) соответственно. Тогда \(UCL\), \(CL\) и \(LCL\) карты Шухарта для \(\mu_w\), когда \(\mu_w\) и \(\sigma_w\) заданы, равны

\begin{equation}
	\begin{aligned}
		UCL &= \mu_w + k\sigma_w \\
		CL &= \mu_w \\
		LCL &= \mu_w - k\sigma_w
	\end{aligned}
	\label{eq:2.1}
\end{equation}

где \(k\) — это «расстояние» верхней (нижней) контрольной границы от \(CL\), выраженное в единицах стандартного отклонения \(W\). Константа \(k\) называется \textbf{константой для нанесения на карту} (charting constant). Эти контрольные границы называются \(k\)-сигма границами, а контрольные карты, построенные в соответствии с этими принципами, называются \(k\)-сигма контрольными картами Шухарта. Популярным выбором для константы \(k\) является 3, что дает 3-сигма контрольную карту Шухарта.

\subsection{Построение контрольных карт}
Статистика для нанесения на карту упоминалась ранее, и при выборе этой статистики обычно используются статистические соображения, такие как несмещенность, минимальная дисперсия, робастность и эффективность. Для мониторинга центральной тенденции или положения процесса, например, часто используется выборочное среднее \(\bar{X}\) или выборочная медиана \(\tilde{X}\). В качестве альтернативы выборочный размах (\(R\)), выборочное стандартное отклонение (\(S\)) и выборочная дисперсия (\(S^2\)) используются для мониторинга разброса, стандартного отклонения или дисперсии процесса соответственно. После выбора статистики для нанесения на карту мы используем ее распределение вероятностей вместе с некоторыми желаемыми критериями эффективности карты, которые будут обсуждаться вскоре, для установления контрольных границ. Многие из классических контрольных карт были разработаны в предположении, что распределение процесса является нормальным. Иногда статистика для нанесения на карту является средним или статистикой, для которой центральная предельная теорема может быть использована для обоснования приблизительной нормальности без явного предположения о нормальности распределения процесса. Когда делается предположение о нормальности или некотором другом распределении (например, экспоненциальном) для процесса, контрольная карта называется \textbf{параметрической} (parametric). Когда такое предположение о распределении не может быть сделано или обосновано, доступны альтернативные контрольные карты, которые называются \textbf{непараметрическими} (nonparametric) или \textbf{свободными от распределения} (distribution-free). Такие карты подробно обсуждаются в Главе 4.

\subsection{Карты по количественным и альтернативным признакам}
Переменные --- это характеристики, которые нас интересуют для мониторинга. Чтобы переменная могла отслеживаться, она должна быть измеримой. Существует два основных типа переменных: количественные и качественные (категориальные). Количественные переменные могут быть непрерывными или дискретными. В SPC термин \textbf{данные по переменным} (variables data) относится к измерениям количественных непрерывных переменных. Примеры включают данные по таким переменным, как длина, ширина, температура, вес, объем и т.д., каждая из которых является непрерывной переменной. В зависимости от параметра, статистикой для нанесения на карту может быть выборочное среднее \(\bar{X}\), выборочная медиана \(\tilde{X}\), выборочный размах (\(R\)), выборочное стандартное отклонение (\(S\)) или выборочная дисперсия (\(S^2\)). Для контрольной карты по количественным признакам статистика для нанесения на карту имеет непрерывное распределение.

И наоборот, термин \textbf{данные по альтернативным признакам} (attributes data) относится к записям или измерениям количественных дискретных переменных. Примеры включают количество ошибок или опечаток, допущенных при заполнении кредитной заявки, или количество врачебных ошибок, допущенных в больнице. Для контрольной карты по альтернативным признакам статистика для нанесения на карту является дискретной и имеет дискретное распределение, например, статистикой для нанесения на карту может быть доля несоответствующих единиц (\(p\)), число несоответствующих единиц (\(np\)), число несоответствий (\(c\)) или среднее число несоответствий на единицу (\(\bar{c}\)).

Контрольные карты по альтернативным признакам имеют преимущество, позволяя быстро обобщать различные аспекты качества продукта, то есть инженер может просто классифицировать продукты как приемлемые или неприемлемые на основе различных критериев качества. Таким образом, карты по альтернативным признакам иногда позволяют избежать необходимости в дорогостоящих, точных устройствах и трудоемких процедурах измерения. Кроме того, этот тип карт, как правило, более понятен менеджерам, не знакомым с процедурами контроля качества; следовательно, он может предоставить более убедительные доказательства (для руководства) проблем с качеством.

Однако в целом, когда все другие соображения равны, данные по переменным и контрольные карты по количественным признакам обычно предпочтительнее и рекомендуются, поскольку данные по переменным и соответствующие карты предоставляют больше информации о процессе. На рисунке \ref{fig-2.6} приведена блок-схема, которая поможет выбрать подходящую контрольную карту Шухарта.

\begin{figure}[!ht]
	\centering
	\includegraphics[width=1.0\linewidth]{"./images/fig-2.6.png"}
	\caption[]{Блок-схема для выбора подходящей контрольной карты для одномерных количественных данных}
	\label{fig-2.6}
\end{figure}

\subsection{Объем выборки или размер подгруппы}
Данные должны быть собраны до того, как контрольная карта может быть внедрена. На практике данные обычно собираются партиями или выборками, каждая из которых обычно содержит \(n\) наблюдений, взятых случайным образом через равные промежутки времени. Партии также называются подгруппами, а такие данные называются сгруппированными в подгруппы. Объем выборки \(n\) обычно представляет собой небольшое число, около пяти, хотя в современной среде существуют приложения, где каждое отдельное наблюдение может быть и проверяется. Такие данные обычно называются \textbf{индивидуальными данными}. В этом случае \(n = 1\). Некоторые случаи, когда отслеживаются индивидуальные значения, включают, например, следующие:

\begin{itemize}
	\item Производство может быть чрезвычайно медленным, и сбор более одного наблюдения нецелесообразен.
	\item Не существует рационального способа разделения данных на подгруппы.
	\item Может быть, выборка размером более одного не оправдана экономически.
\end{itemize}

Мы рассмотрим случай индивидуальных наблюдений позже в разделе о контрольных картах для индивидуальных наблюдений.

\subsection{Рациональное разделение на подгруппы}
Рациональная подгруппа — это выборка или подгруппа, взятая из процесса таким образом, что вариация между измерениями или наблюдениями внутри выборки обусловлена только присущей или естественной изменчивостью процесса, и никакая дополнительная вариация из-за каких-либо особых причин не присутствует внутри выборочных измерений. Рациональные подгруппы — это те подгруппы, которые содержат или включают только общие/случайные причины. Вариация между показаниями должна быть обусловлена естественной вариацией процесса. Не должно быть возможности для того, чтобы назначаемые/особые причины добавлялись к этой вариации.

\noindent\textbf{Пример 2.1. Рациональное разделение на подгруппы}

Предположим, у нас есть два оператора, работающих на одном и том же станке в разное время, и мы хотим выявить различия между операторами, используя выборки объёмом 3.

\begin{center}
	\(\bullet\) Оператор 1 \\
	\(\ast\) Оператор 2
\end{center}

Тогда правильный способ рационального разбиения на подгруппы будет, например, таким:

\begin{center}
	\begin{tikzpicture}[x=0.35cm, y=0.5cm]
		% Символы (Оператор 1 - точки, Оператор 2 - звездочки)
		\node at (0.5, 0) {\(\bullet\)};
		\node at (1.5, 0) {\(\bullet\)};
		\node at (2.5, 0) {\(\bullet\)};
		\node at (3.5, 0) {\(\bullet\)};
		\node at (4.5, 0) {\(\bullet\)};
		
		\node at (6.5, 0) {\(\ast\)};
		
		\node at (8.5, 0) {\(\ast\)};
		\node at (9.5, 0) {\(\ast\)};
		\node at (10.5, 0) {\(\ast\)};
		
		\node at (12.5, 0) {\(\ast\)};
		
		\node at (14.5, 0) {\(\bullet\)};
		\node at (15.5, 0) {\(\bullet\)};
		\node at (16.5, 0) {\(\bullet\)};
		\node at (17.5, 0) {\(\bullet\)};
		\node at (18.5, 0) {\(\bullet\)};
		
		\node at (20.5, 0) {\(\ast\)};
		\node at (21.5, 0) {\(\ast\)};
		
		\node at (23.5, 0) {\(\ast\)};
		\node at (24.5, 0) {\(\ast\)};
		\node at (25.5, 0) {\(\ast\)};
		
		% Фигурные скобки для правильного разбиения
		\draw[decorate, decoration={brace, amplitude=5pt, mirror}] 
		(2, -0.2) -- (5, -0.2);
		\draw[decorate, decoration={brace, amplitude=5pt, mirror}] 
		(8, -0.2) -- (11, -0.2);
		\draw[decorate, decoration={brace, amplitude=5pt, mirror}] 
		(14, -0.2) -- (17, -0.2);
		\draw[decorate, decoration={brace, amplitude=5pt, mirror}] 
		(23, -0.2) -- (26, -0.2);
	\end{tikzpicture}
\end{center}

В этом случае подгруппы выбираются так, чтобы изолировать работу каждого оператора. Неправильным способом рационального разделения на подгруппы будет, например:

\begin{center}
	\begin{tikzpicture}[x=0.35cm, y=0.5cm]
		% Символы (Смешанные группы)
		\node at (0.5, 0) {\(\bullet\)};
		\node at (1.5, 0) {\(\bullet\)};
		\node at (2.5, 0) {\(\bullet\)};
		\node at (3.5, 0) {\(\bullet\)};
		\node at (4.5, 0) {\(\bullet\)};
		
		\node at (6.5, 0) {\(\ast\)};
		\node at (7.5, 0) {\(\ast\)};
		\node at (8.5, 0) {\(\ast\)};
		\node at (9.5, 0) {\(\ast\)};
		\node at (10.5, 0) {\(\ast\)};
		
		\node at (12.5, 0) {\(\bullet\)};
		\node at (13.5, 0) {\(\bullet\)};
		\node at (14.5, 0) {\(\bullet\)};
		\node at (15.5, 0) {\(\bullet\)};
		\node at (16.5, 0) {\(\bullet\)};
		
		% Фигурные скобки для неправильного разбиения
		\draw[decorate, decoration={brace, amplitude=5pt, mirror}] 
		(4, -0.2) -- (8, -0.2);
		\draw[decorate, decoration={brace, amplitude=5pt, mirror}] 
		(10, -0.2) -- (14, -0.2);
	\end{tikzpicture}
\end{center}

Действуя таким образом, любые различия между операторами будут эффективно замаскированы.

Рисунок 2.7 представляет процесс, который состоит из пяти станков, объединяющих свою продукцию в один «бункер». Предположим, что из «бункера» берётся случайная выборка для формирования подгруппы объёмом 5. Мы могли бы получить, например,\(\circ\blacktriangle\square\circ\bullet\) или \(\square\square\bullet\blacksquare\blacktriangle\). Такие подгруппы не смогут обнаружить различия между станками. Это называется \textbf{эффектом маскировки} (masking effect), поскольку различия между станками будут эффективно скрыты. Одно из решений — брать подгруппу объёмом 5 с первого станка, затем подгруппу объёмом 5 со второго станка и так далее. Выбор подгрупп таким образом гарантирует, что контрольная карта способна обнаруживать различия между станками.

\begin{figure}[htbp]
	\centering
	\includegraphics[width=0.8\textwidth]{"./images/fig-2.7.png"}
	\caption{Пример, в котором различия между процессами оказываются эффективно замаскированы.}
	\label{fig-2.7}
\end{figure}

\subsection{Непараметрические или свободные от распределения методы}
С исторической точки зрения, большая часть ранних работ в статистике была основана на предположении о нормальности лежащего в основе распределения. Работа в SPC, по большому счету, не является исключением. Хотя большинство традиционных контрольных карт по количественным признакам были параметрическими, то есть они основаны на некотором параметрическом предположении о модели распределения процесса, таком как нормальность, существует множество ситуаций на практике, когда такое предположение нежизнеспособно или неоправданно и может, фактически, быть контрпродуктивным. В таких ситуациях может быть предпочтительнее использовать альтернативный набор контрольных карт, называемых \textbf{непараметрическими} или \textbf{свободными от распределения}. Как и в Gibbons and Chakraborti (2010), термин «непараметрический» не предназначен для того, чтобы подразумевать, что параметры не задействованы; на самом деле, совсем наоборот. Хотя термин «свободный от распределения» кажется лучшим описанием того, что мы ожидаем от этих карт, то есть они остаются «валидными» для всех непрерывных распределений, «непараметрический» — это, пожалуй, термин, который используется чаще. Под валидностью мы подразумеваем, что характеристики IC-эффективности этих карт остаются неизменными для всех непрерывных распределений, независимо от их формы. Таким образом, непараметрические контрольные карты имеют гарантированную IC-робастность, что является их главным преимуществом перед параметрическими картами. Более того, теперь было продемонстрировано, что непараметрические карты часто имеют конкурентоспособную OOC-эффективность по сравнению с их параметрическими аналогами, и, таким образом, их использование на практике кажется беспроигрышным вариантом.

В статистической литературе в настоящее время существует довольно обширная коллекция непараметрических тестов гипотез и доверительных интервалов, и было показано, что эти методы хорошо работают по сравнению с их аналогами из нормальной теории. Примечательно, что даже когда лежащее в основе распределение является нормальным, асимптотическая относительная эффективность некоторых непараметрических методов по отношению к соответствующим (оптимальным) методам нормальной теории может достигать 0,955 (см., например, Gibbons and Chakraborti, 2010, с. 508). Для некоторых других распределений с тяжелыми хвостами и скошенных распределений эффективность может быть 1,0 или даже выше. Можно возразить, что непараметрические методы будут «менее эффективными», чем их параметрические аналоги, когда имеется полное знание распределения процесса, для которого этот параметрический метод был специально разработан. Однако реальность такова, что такая информация редко, если вообще когда-либо, доступна на практике.

Параметрические контрольные карты используются наиболее часто, поскольку они разрабатывались с самого начала SPC, примерно с 1940-х годов. Они доступны во многих программных пакетах. Однако с практической точки зрения кажется естественным разрабатывать и использовать непараметрические методы, включая контрольные карты, в приложениях SPC. Практикам в области качества было бы полезно иметь эти методы в своем арсенале.

\begin{enumerate}
	\item Это легко реализовать, то есть простота.
	\item Нет необходимости предполагать конкретное параметрическое распределение для лежащего в основе процесса.
	\item Распределение длины серии в состоянии IC одинаково для всех непрерывных распределений процесса.
	\item Это более робастно и устойчиво к выбросам.
	\item Это может быть более эффективным при обнаружении изменений, когда истинное распределение заметно отличается от нормального, особенно с более тяжелыми хвостами.
	\item Оценка дисперсии для построения карт для параметра положения может не всегда быть необходимой.
\end{enumerate}

С технической точки зрения, большинство непараметрических процедур, включая контрольные карты, требуют, чтобы генеральная совокупность была непрерывной, чтобы быть свободными от распределения. Таким образом, в контексте SPC мы рассматриваем так называемые непараметрические контрольные карты по количественным признакам. Хотя непараметрические карты предлагают много преимуществ, у них есть и критики. Некоторые из воспринимаемых недостатков непараметрических контрольных карт, наряду с нашими ответами (в скобках), следующие:

\begin{enumerate}
	\item Они будут «менее эффективными», чем их параметрические аналоги, когда имеется полное знание распределения процесса, для которого этот параметрический метод был специально разработан (но такое полное знание редко, если вообще когда-либо, доступно).
	\item Обычно требуются специальные таблицы, когда размеры выборки малы (компьютерное программное обеспечение может быть написано и сделано доступным, чтобы преодолеть это).
	\item Непараметрические методы не так хорошо известны среди исследователей и практиков в области качества (это может быть исправлено с помощью большего обучения и освещения).
\end{enumerate}

Хотя непараметрические контрольные карты составляют основную тему этой книги, чтобы подготовить почву, мы сначала обсудим некоторые детали традиционных параметрических карт в Главе 3. После этого, в Главе 4, мы обсудим одномерные непараметрические контрольные карты, которые предназначены для отслеживания положения непрерывного процесса; для масштаба доступно очень мало карт, и в этой области требуется больше работы. Область многомерных контрольных карт интересна, и объем литературы по непараметрическим многомерным контрольным картам растет. Однако мы не фокусируемся на них в этой книге; читатель отсылается к недавней книге Qiu (2014), в которой есть хорошее обсуждение некоторых параметрических и непараметрических контрольных карт.

Основной момент таков: если можно предположить лежащую в основе модель, и эта модель окажется правильной, параметрическая статистическая процедура, такая как параметрическая контрольная карта, была бы более эффективной, чем непараметрическая контрольная карта, которая не делает предположений о модели. Однако реальность такова, что на практике такое предположение о модели часто может быть недоступным или неоправданным, и, как таковой, непараметрическая контрольная карта может часто быть единственной жизнеспособной альтернативой для использования.

\subsection{Мониторинг положения процесса и/или масштаба процесса}
Две важные классические проблемы в обычном SPC — это мониторинг среднего значения процесса и/или стандартного отклонения процесса. В непараметрической постановке мы рассматриваем, в более общем смысле, мониторинг центра или параметра положения и/или параметра масштаба процесса. Параметр положения представляет типичное значение и может быть средним или некоторым процентилем, например, медианой распределения. Последнее особенно привлекательно, когда ожидается, что лежащее в основе распределение будет скошенным. Пусть \(F(x)\) обозначает неизвестную функцию распределения (cdf) отслеживаемой непрерывной переменной \(X\). По аналогии с параметрическим, главным образом нормальным, случаем, предполагается, что \(F\) следует либо:

\begin{enumerate}
	\item модели положения, с cdf \(F(x - \theta)\), где \(x \in (-\infty, \infty)\) и \(\theta \in (-\infty, \infty)\) — параметр положения; или
	\item модели масштаба, с cdf \(F\left(\frac{x}{\tau}\right)\), где \(x \in (-\infty, \infty)\) и \(\tau > 0\) — параметр масштаба; или
	\item модели положения-масштаба с cdf \(F\left(\frac{x - \theta}{\tau}\right)\), где \(x \in (-\infty, \infty)\), \(\theta \in (-\infty, \infty)\) и \(\tau > 0\) — параметры положения и масштаба соответственно.
\end{enumerate}

Таким образом, в непараметрической постановке цель состоит в мониторинге \(\theta\) или \(\tau\), или обоих, в рамках этих модельных предположений на основе случайных выборок или подгрупп размером \(n \geq 1\) данных, обычно берущихся через равные промежутки времени, или с \(n = 1\), что соответствует случаю индивидуальных измерений. Чтобы подчеркнуть аналогию в обычных задачах параметрического контроля, предполагается, что \(F\) является cdf (\(\Phi\)) стандартного нормального распределения, тогда как в непараметрической постановке, для данных по переменным, \(F\) — это некоторая неизвестная непрерывная cdf. Хотя модель положения-масштаба кажется естественной моделью для рассмотрения, параллельной случаю нормальной теории с неизвестными средним и дисперсией, большая часть того, что в настоящее время доступно в литературе по непараметрическому статистическому контролю процессов (NSPC), в основном имеет дело с моделью положения. Как отмечалось ранее, это будет основным фокусом в этой книге.

\subsection{Случай K и случай U}

Обратите внимание, что для параметрических контрольных карт делается предположение о распределении для лежащего в основе процесса, такое как нормальность, но параметры распределения могут быть известны/заданы или неизвестны. Когда лежащие в основе параметры известны или заданы, это называется «случаем известных стандартов» и обозначается как \textbf{Случай K} (Case K) для удобства. Напротив, если параметры неизвестны и должны быть оценены, эта ситуация называется «случаем неизвестных стандартов» и обозначается как \textbf{Случай U} (Case U). Обычно оценка параметров выполняется в фазе I с использованием эталонных данных. Для параметрических карт оба случая изучаются в Главе 3. Для непараметрических карт, хотя мы не делаем конкретных модельных предположений, таких как нормальность, у нас все еще есть Случай K и Случай U, как обсуждается в Главе 4. Мы сначала рассмотрим Случай K, а затем случай, более типичный для промышленности, которым является Случай U. Мы начнем с обсуждения параметрических контрольных карт по количественным признакам в Главе 3.

\subsection{Контрольные карты и проверка гипотез}

Существует сходство между контрольными картами и традиционными инструментами вывода, такими как проверка гипотез и доверительные интервалы. Читатель должен помнить, что при использовании карты Шухарта процесс объявляется IC, когда статистика для нанесения на карту попадает между контрольными границами, тогда как процесс объявляется OOC, когда статистика для нанесения на карту попадает на или за пределы любой контрольной границы. Таким образом, как и при проверке гипотез, существует решение «да/нет» на основе значения статистики для нанесения на карту и областей принятия решений.

Предположим, что статистика для нанесения на карту \(T\) со стандартным отклонением \(\sigma_T\) является несмещенной точечной оценкой параметра процесса \(\theta\). То есть \(E(T) = \theta\) и \(VAR(T) = \sigma_T^2\). Тогда существует связь между использованием контрольной карты Шухарта на основе статистики для нанесения на карту \(T\) и проверкой гипотезы вида

\[
H_0: \theta = \theta_0 \quad \text{против} \quad H_1: \theta \neq \theta_0.
\]

Например, если \(T\) — выборочное среднее (т.е. \(\overline{X}\)), и мы хотим проверить гипотезу о том, что среднее процесса \(\mu\) равно некоторому заданному значению \(\mu_0\), проблема проверки гипотез имеет вид

\[
H_0: \mu = \mu_0 \quad \text{против} \quad H_1: \mu \neq \mu_0
\]

где \(H_0\) — нулевая гипотеза, а \(H_1\) — альтернативная гипотеза.

Обычно \(H_0\) отвергается тогда и только тогда, когда \(\frac{\sqrt{n}|\overline{X} - \mu_0|}{\sigma_0} \geq k\), или, альтернативно, тогда и только тогда, когда

\begin{equation}
	\overline{X} \geq \mu_0 + k\frac{\sigma_0}{\sqrt{n}}
	\label{eq:2.2}
\end{equation}

или

\begin{equation}
	\overline{X} \leq \mu_0 - k\frac{\sigma_0}{\sqrt{n}}.
	\label{eq:2.3}
\end{equation}

Обратите внимание, что правые части Уравнений 2.2 и 2.3 фактически являются \(k\)-сигма контрольными границами из Уравнения 2.1. Таким образом, если \(\overline{X}\) отображается между двумя контрольными границами, мы заключаем, что среднее процесса находится в состоянии статистической управляемости, что равносильно утверждению, что мы не отвергаем нулевую гипотезу \(H_0: \mu = \mu_0\). И наоборот, если \(\overline{X}\) отображается на или за пределами любой контрольной границы, мы заключаем, что среднее процесса не находится в состоянии статистической управляемости, что, в свою очередь, означает, что та же нулевая гипотеза \(H_0: \mu = \mu_0\) отвергается в пользу альтернативной гипотезы \(H_1: \mu \neq \mu_0\). Таким образом, две OOC-области контрольной карты Шухарта можно рассматривать как графическое отображение области отклонения или критической области для проверки гипотезы вида \(H_0: \theta = \theta_0\) против \(H_1: \theta \neq \theta_0\). Следовательно, контрольную карту можно рассматривать как проверку нулевой гипотезы \(H_0: \theta = \theta_0\) повторно в ряде моментов времени. В литературе существуют различные точки зрения на эти вопросы и споры по ним, см., например, статью Woodall (2000) вместе с вкладами различных участников дискуссии.

Структура проверки гипотез может быть полезна при изучении эффективности контрольной карты, например, как и в контексте проверки гипотез, здесь мы также можем думать об ошибках и вероятностях ошибок. Ошибка I рода — это событие, когда процесс объявляется OOC, хотя на самом деле процесс IC. Аналогично, ошибка II рода — это событие, когда процесс объявляется IC, хотя на самом деле процесс не IC. Это представлено в Таблице 2.1.

Соответствующие вероятности этих двух ошибок (или событий) равны

\[
\alpha = P(\text{Ошибка I рода}) = P(\text{Объявить процесс OOC} \mid \text{процесс на самом деле IC})
\]

и

\[
\beta = P(\text{Ошибка II рода}) = P(\text{Объявить процесс IC} \mid \text{процесс на самом деле OOC})
\]

соответственно.

\begin{table}[htbp]
	\centering
	\caption{Представление ошибок I и II рода.}
	\label{tab:2.1}
	\begin{tabular}{lll}
		\toprule
		& \multicolumn{2}{c}{Состояние процесса признается} \\
		\cmidrule{2-3}
		Состояние процесса & Управляемым (IC) & Неуправляемым (OOC) \\
		\midrule
		Управляемое (IC) & Верное решение & Ошибка I рода \\
		Неуправляемое (OOC) & Ошибка II рода & Верное решение \\
		\bottomrule
	\end{tabular}
\end{table}

В приложениях SPC ошибка I рода обычно называется \textbf{ложной тревогой}, а связанная с ней вероятность называется \textbf{частотой ложных тревог} (FAR). FAR подробно обсуждается в Разделе 2.1.12. Поскольку вероятность ошибки I рода более серьезна при проверке гипотез, FAR считается более серьезной в SPC, и обычно стараются контролировать FAR на довольно малом уровне, например, значение FAR 0,0027 часто используется в промышленности, что аналогично использованию значения \(\alpha\) 0,05 в традиционном статистическом анализе.

Вероятность ошибки II рода называется \textbf{\(\beta\)-риском}. Две вероятности ошибок обратно связаны друг с другом. Если FAR увеличивается, \(\beta\)-риск уменьшается, и наоборот. Возникает вопрос: почему \(\alpha\) и \(\beta\) важны при выборе контрольных границ?

Если мы расширим контрольные границы (выберем большее значение \(k\)), это снизит FAR, что уменьшит риск ошибки I рода, что уменьшит риск того, что точка попадет на или за пределы контрольных границ, когда назначаемые причины отсутствуют. Аналогично, если мы расширим контрольные границы (выберем большее значение \(k\)), это увеличит риск ошибки II рода, что увеличит риск того, что точка попадет между контрольными границами, когда процесс действительно OOC.

\subsection{Общие шаги при проектировании контрольной карты}
В предыдущих разделах обсуждались многие вопросы, касающиеся контрольных карт, такие как: следует ли отслеживать положение и/или масштаб, следует ли использовать карты по альтернативным или количественным признакам, выбор объема выборки, рациональное разделение на подгруппы и т.д. Здесь мы суммируем общие шаги, которым нужно следовать при проектировании (любой) контрольной карты, прежде чем вводить различные контрольные карты.

Пусть \(X_{i1}, X_{i2}, \ldots, X_{in}\) обозначают случайную выборку (измерения некоторой количественной характеристики качества) объема \(n \geq 1\), взятую из процесса в момент времени \(i = 1, 2, 3, \ldots\)

\textbf{Шаг 1.} Определить или выбрать лежащее в основе распределение процесса (скажем, \(F(x)\)) или считать его неизвестным, но непрерывным. Это приведет к рассмотрению параметрической или непараметрической контрольной карты. Определить структуру корреляции (i.i.d. или не-i.i.d. наблюдения). Пример: предположим, \(X_{ij} \sim \text{i.i.d. } F(x)\), где \(F(x)\) — это cdf нормального распределения или \(F(x)\) — неизвестная непрерывная cdf.

\textbf{Шаг 2.}
\begin{itemize}
	\item Решить, какую(ие) характеристику(и) процесса (параметр(ы)) нужно отслеживать (например, положение и/или разброс).
	\item Решить, какую(ие) выборочную(ые) статистику(и) использовать (например, \(\overline{X}\), \(\tilde{X}\), \(R\), \(S\), \(S^2\) и т.д.). В общем случае обозначим эту статистику через \(T_j = h(x_{j1}, \ldots, x_{jn})\), где \(h\) — некоторая функция.
\end{itemize}

\textbf{Шаг 3.} Решить, какой тип контрольной карты использовать (например, Шухарта, CUSUM или EWMA) или любое из их улучшений. Это решение будет зависеть от того, какие сдвиги представляют интерес, наряду с другими практическими соображениями. Пример: предполагая нормальность, использовать CUSUM-карту для \(T_j = \overline{X}_j\) для обнаружения малого сдвига среднего.

\textbf{Шаг 4.} Получить распределение длины серии контрольной карты и, используя его, найти и использовать константу для нанесения на карту так, чтобы получить желаемую эффективность. Желаемая эффективность включает критерий, связанный с распределением длины серии, такой как средняя длина серии, обозначаемая \(ARL\). Различные критерии эффективности, такие как \(ARL\), определяются и обсуждаются ниже.

\subsection{Показатели эффективности контрольной карты}

Поскольку контрольная карта используется для мониторинга процесса, а качество выходного продукта процесса жизненно важно, эффективность любой контрольной карты является важным соображением. Различные характеристики распределения длины серии обычно используются для оценки эффективности карты. \textbf{Длина серии} (run-length) --- это случайная величина, обозначающая количество точек (статистик для нанесения на карту), которые должны быть нанесены на график (или количество подгрупп, которые должны быть собраны, каждая из которых дает, скажем, одну статистику для нанесения на карту), прежде чем карта подаст первый OOC-сигнал (или будет наблюдаться первое сигнальное событие). Длина серии обозначается через \(N\). Поскольку длина серии — это считающая случайная величина, она дискретна, имеет функцию вероятности (pmf), и мы изучаем различные атрибуты этой pmf для оценки эффективности карты. Например, мы вычисляем ее среднее или ожидаемое значение, так называемую \textbf{среднюю длину серии} (обозначаемую \(ARL\)), чтобы понять среднюю эффективность, ее стандартное отклонение, которое является стандартным отклонением распределения длины серии (обозначаемым \(SDRL\)), чтобы понять изменчивость эффективности. Можно видеть, что pmf длины серии имеет правостороннюю асимметрию, и поэтому мы также рассматриваем ее процентили, такие как медиана, так называемая \textbf{медианная длина серии} (обозначаемая \(MRL\)), для более полного понимания эффективности карты. \(ARL\), возможно, является наиболее известным показателем эффективности карты на практике, но \(MRL\) в последнее время набирает обороты, в то время как были сделаны и другие предложения. Мы обсудим некоторые из этих вопросов более подробно позже.

Для дальнейшего обсуждения важно сначала подчеркнуть тот факт, что SPC обычно состоит из двух частей (или фаз, или этапов), а именно \textbf{Фазы I} (или Этапа I) и \textbf{Фазы II} (или Этапа II), в которых цели мониторинга различны, так что контрольные карты строятся с использованием разных критериев эффективности. С этого момента мы будем ссылаться на фазы отдельно, чтобы сделать вопросы более ясными. Первая фаза, известная как \textbf{Фаза I} или ретроспективная фаза, в основном носит исследовательский характер и обычно используется для установления контроля или стабильности процесса на основе анализа исторических или ретроспективных данных. Обычно эти данные тщательно анализируются с использованием описательных и инференциальных статистических методов, включая контрольные карты. Затем оцениваются неизвестные параметры, проверяются предположения о распределении и рассчитываются контрольные границы. Контрольные карты Фазы I рассчитываются с использованием несколько иного критерия, о котором будет сказано больше позже. Затем рассчитывается и наносится на контрольную карту статистика для нанесения на карту из каждой доступной выборки. Если все статистики для нанесения на карту находятся в пределах контрольных границ и никакие систематические паттерны не очевидны, процесс объявляется IC. Если же одна или несколько контрольных статистик попадают на или за пределы контрольных границ, обычно исследуются эти выборки на предмет назначаемых причин. Если таковые присутствуют, эти выборки могут быть отброшены, а контрольные границы пересчитаны с использованием оставшихся выборок. Затем статистики для нанесения на карту снова наносятся на пересмотренную контрольную карту. Этот процесс повторяется до тех пор, пока мы не будем удовлетворены тем, что процесс находится в состоянии IC, после чего имеющиеся данные называются \textbf{эталонными данными}; они образуют основу, на которой оцениваются любые неизвестные параметры и рассчитываются контрольные границы с использованием этих оценок. Следующая часть SPC, а именно будущий мониторинг процесса, называется \textbf{проспективной фазой} или \textbf{Фазой II}. Таким образом, анализ Фазы I обычно включает вопросы, связанные с проектированием, оценкой (при необходимости) и установлением контрольных границ итеративным образом, тогда как Фаза II включает будущий мониторинг процесса с использованием оценок параметров и установление контрольных границ с информацией и другими предположениями, разработанными в ходе анализа Фазы I.

Понимание статистической эффективности контрольной карты важно в обеих фазах. В Фазе I обычно исследуется \textbf{вероятность ложной тревоги} (FAP) или вероятность хотя бы одной ложной тревоги (сигнал, когда процесс IC). Однако в Фазе II некоторые атрибуты распределения длины серии, такие как \(ARL\), \(SDRL\) или \(MRL\), используются для измерения и описания эффективности карты. Сначала мы обсудим показатели эффективности контрольной карты для Фазы I (т.е. FAP), после чего обсудим показатели эффективности контрольной карты для Фазы II.

\subsubsection{Вероятность ложной тревоги (FAP)}
Без потери общности предположим, что \(m\) обозначает «окончательное» количество эталонных выборок в конце анализа Фазы I, когда процесс считается IC. Таким образом, предполагается, что набор эталонных данных содержит \(N = mn\) наблюдений. Ситуация OOC указывается, когда статистика для нанесения на карту попадает на или за пределы любой из контрольных границ. Как отмечалось ранее, это важное событие часто называется \textbf{сигналом} или \textbf{сигнальным событием}. Удобно рассмотреть дополнительное событие, то есть когда подгруппа не подает сигнал, и оно называется \textbf{несигнальным событием}. Таким образом, для \(i\)-й подгруппы \(E_i = \{L\hat{C}L < C_i < U\hat{C}L\}\), \(i = 1, 2, \ldots, m\) обозначает несигнальное событие. Обратите внимание, что \(\hat{L}\hat{C}L\) и \(\hat{U}\hat{C}L\) обозначают оцененные контрольные границы Фазы I, рассчитанные с использованием эталонных данных, а \(C_i\) обозначает статистику для нанесения на карту из \(i\)-й эталонной выборки. Как обсуждалось ранее, FAP обычно является рекомендуемым критерием проектирования карты, принятым в Фазе I (см. Chakraborti, Human, and Graham 2009). FAP можно выразить следующим образом:

\begin{equation}
	\begin{aligned}
		FAP &= P(\text{Хотя бы одна ложная тревога из } m \text{ подгрупп}) \\
		&= 1 - P(\text{Ни одного сигнала среди } m \text{ подгрупп} \mid \text{IC}) \\
		&= 1 - P(E_1, E_2, \ldots, E_m \mid \text{IC}) \\
		&= 1 - P\left(\bigcap_{i=1}^m \{L\hat{C}L < C_i < U\hat{C}L\} \mid \text{IC}\right).
	\end{aligned}
	\label{eq:2.4}
\end{equation}

Последний шаг Уравнения 2.4 равен

\[
1 - \int_{L\hat{C}L}^{U\hat{C}L} \int_{L\hat{C}L}^{U\hat{C}L} \cdots \int_{L\hat{C}L}^{U\hat{C}L} f_{C_1, C_2, \ldots, C_m}(c_1, c_2, \ldots, c_m) dc_1 \cdots dc_m
\]

если совместное распределение \(C_i\) непрерывно, и где \(f_{C_1, C_2, \ldots, C_m}(c_1, c_2, \ldots, c_m)\) обозначает совместную функцию плотности вероятности статистик для нанесения на карту \(C_1, C_2, \ldots, C_m\), когда процесс IC. Обратите внимание, что совместное распределение используется для расчета FAP, поскольку \(m\) сигнальных событий в этом случае являются зависимыми, так как каждая статистика для нанесения на карту сравнивается с одними и теми же границами, которые построены из одного и того же набора из \(m\) эталонных выборок.

Аналогично, если \(C_i\) дискретны, то \(P(\cap_{i=1}^m \{L\hat{C}L < C_i < U\hat{C}L\} \mid \text{IC})\) можно переписать как \(P(\cap_{i=1}^m \{L\hat{C}L + 1 \leq C_i \leq U\hat{C}L - 1\} \mid \text{IC})\), и, следовательно, Уравнение 2.4 равно

\[
= 1 - \sum_{c_1=L\hat{C}L+1}^{U\hat{C}L-1} \sum_{c_2=L\hat{C}L+1}^{U\hat{C}L-1} \cdots \sum_{c_m=L\hat{C}L+1}^{U\hat{C}L-1} f_{C_1, C_2, \ldots, C_m}(c_1, c_2, \ldots, c_m)
\]

где \(f_{C_1, C_2, \ldots, C_m}(c_1, c_2, \ldots, c_m)\) обозначает совместную pmf статистик для нанесения на карту \(C_1, C_2, \ldots, C_m\), когда процесс IC.

Таким образом, расчет FAP требует знания совместного распределения статистик для нанесения на карту, когда процесс IC. Вывод этого распределения и определение контрольных границ (и связанных с ними констант для нанесения на карту; которые различаются в зависимости от типа рассматриваемой контрольной карты) представляют собой некоторые важные практические проблемы, особенно в контексте SPC, где количество сравнений (\(m\)) может быть как малым, например 25, так и большим, например 300 или более.

FAP и аналогично вероятность сигнала, то есть вероятность хотя бы одного сигнала от \(m\) подгрупп, являются единственными показателями эффективности, обычно рассматриваемыми для Фазы I. Далее мы обсудим показатели эффективности для Фазы II.

\subsubsection{Частота ложных тревог (FAR)}

Ошибка I рода — это событие, когда процесс объявляется OOC, хотя на самом деле процесс IC. Это концепция, используемая в первую очередь в проверке гипотез. В приложениях SPC ошибка I рода обычно называется \textbf{ложной тревогой} (false
alarm), а связанная с ней вероятность называется \textbf{частотой ложных тревог} (false alarm
rate), что сокращается как FAR. Значение \(FAR = 0,0027\) является типичным отраслевым стандартом, но в зависимости от конкретного применения используются и другие значения, такие как 0,01 или 0,001.

\subsubsection{Средняя длина серии (ARL)}

Как уже отмечалось, эффективность контрольной карты изучается через распределение длины серии, и наиболее популярным показателем эффективности является ожидаемое значение распределения длины серии, так называемая \textbf{средняя длина серии} (average
run-length -- ARL). ARL контрольной карты --- это ожидаемое количество статистик для нанесения на карту, которые должны быть нанесены на график (подгрупп, которые должны быть собраны), прежде чем контрольная карта подаст сигнал в первый раз. Когда процесс IC, ожидаемое количество подгрупп, которые должны быть собраны, прежде чем контрольная карта подаст сигнал в первый раз, называется \textbf{IC средней длиной серии} и обозначается \(ARL_{\text{IC}}\). И наоборот, \textbf{OOC средняя длина серии}, обозначаемая \(ARL_{\text{OOC}}\), — это ожидаемое количество выборок, которые должны быть собраны, прежде чем контрольная карта подаст сигнал в первый раз, когда процесс вышел из-под контроля. Интуитивно понятно, что для эффективной контрольной карты \(ARL_{\text{IC}}\) должна быть большой, а \(ARL_{\text{OOC}}\) — маленькой. Следовательно, при сравнении эффективности двух или более карт \(ARL_{\text{IC}}\) карт фиксируется на приемлемо высоком уровне, и карта с наименьшей \(ARL_{\text{OOC}}\) для некоторого заданного сдвига (\(\delta\)) параметра процесса работает лучше всего. Обратите внимание, что эта концепция аналогична фиксации вероятности ошибки I рода некоторых тестов гипотез на номинальном значении, таком как 0,05, и выбору среди них теста с наибольшей мощностью при некотором значении параметра в рамках альтернативной гипотезы. На Рисунке 2.9 мы построили две кривые \(ARL\) в зависимости от сдвигов разного размера, соответствующих двум разным контрольным картам.

Из Рисунка \ref{fig-2.8} видно, что значения \(ARL_{\text{IC}}\) двух карт зафиксированы на номинальном уровне \(ARL_{\text{IC}} = 370\), что указывает на то, что, когда процесс IC, первый сигнал будет наблюдаться в среднем на каждой 370-й нанесенной точке. Номинальное значение 370 является отраслевым стандартом. При этом фиксация значений \(ARL_{\text{IC}}\) делается для того, чтобы мы могли найти более мощную из двух карт, изучив значения \(ARL_{\text{OOC}}\). Наиболее мощная карта — это та, которая дает наименьшую \(ARL_{\text{OOC}}\) для любого заданного сдвига. 

\begin{figure}[!ht]
	\centering
	\includegraphics[width=0.7\linewidth]{"./images/fig-2.8.png"}
	\caption[]{Кривые средней длины серии (ARL) для двух различных контрольных карт в зависимости от величины сдвига}
	\label{fig-2.8}
\end{figure}

Из Рисунка \ref{fig-2.8} видно, что контрольная карта, обозначенная пунктирной кривой, является лучшим выбором, так как она работает лучше равномерно для всех значений сдвига.

\subsubsection{Стандартное отклонение длины серии (SDRL)}

Другим показателем эффективности карты является стандартное отклонение распределения длины серии, обозначаемое как SDRL. Очевидно, что для желательной контрольной карты мы хотим, чтобы SDRL был малым, а ARL был равен или близок к номинальному значению \(ARL_{\text{IC}}\).

\subsubsection{Процентили длины серии}

Хотя \(ARL\) является наиболее часто используемым показателем эффективности карты, отмечая, что распределение длины серии значительно асимметрично вправо, многие исследователи рекомендовали изучать ряд процентилей, включая 5-й, 25-й, 50-й, 75-й и 95-й процентили, чтобы лучше охарактеризовать распределение длины серии; это обсуждается далее.

Оценка процентилей распределения длины серии важна для лучшего понимания эффективности карты, поскольку распределение длины серии значительно асимметрично вправо. Вспомните обсуждение медианы, которая является 50-м процентилем, в Главе 1. Мы отметили, что в общем случае процентили должны быть определены так, чтобы они были уникальными, что может быть проблемой для дискретных распределений, поскольку cdf имеет скачки в точках массы. С этой целью, поскольку случайная величина длины серии \(N\) дискретна, мы начнем с определения \(100 \cdot q\)-го процентиля (\(0 < q < 1\)) как наименьшего положительного целого числа \(t\), такого что \(P(N \leq t) \geq q\). Другими словами, \(100 \cdot q\)-й процентиль распределения — это наименьшее положительное целое число \(t\), такое что cdf \(N\) в \(t\) по крайней мере равна \(q\). Наиболее популярными процентилями для исследования являются 5-й, 25-й, 50-й, 75-й и 95-й процентили, которые охватывают распределение и включают первый, второй и третий квартили. Второй квартиль или 50-й процентиль — это медиана, а медианная длина серии обозначается \(MRL\). Идея рассмотрения процентилей в SPC восходит к Barnard (1959). Многие авторы, включая Chakraborti (2007) и Khoo et al. (2011), выступали за использование процентилей, таких как медиана, для оценки эффективности карты. Предположим, что \(MRL\) используется как мера типичной эффективности карты. \(MRL\) может быть предпочтительной мерой с точки зрения робастности, поскольку, как мы отмечали ранее, распределение длины серии значительно асимметрично вправо. Тогда эффективная контрольная карта должна иметь большую \(MRL_{\text{IC}}\) и малую \(MRL_{\text{OOC}}\). Это аналогично по духу подходу при использовании \(ARL\) в качестве показателя эффективности. При сравнении эффективности карт среди карт с приемлемо высокой номинальной \(MRL_{\text{IC}}\), предпочтительной является контрольная карта с наименьшей \(MRL_{\text{OOC}}\) при некотором заданном сдвиге (\(\delta\)) параметра процесса. На Рисунке 2.10 мы построили две кривые \(MRL\), соответствующие двум разным контрольным картам, в зависимости от сдвигов разного размера.

Из Рисунка \ref{fig-2.9} видно, что значения \(MRL_{\text{IC}}\) карт зафиксированы на номинальном уровне \(MRL_{\text{IC}} = 350\), что указывает на то, что первая ложная тревога будет наблюдаться в пределах первых 350 нанесенных точек, по крайней мере, в 50\% случаев. Это интерпретация, которой нет у \(ARL_{\text{IC}}\). Фиксация \(MRL_{\text{IC}}\) на уровне 350 является отраслевым стандартом. При этом, поскольку обе карты имеют одинаковое значение \(MRL_{\text{IC}}\), мы можем найти более эффективную из двух карт, изучив значения \(MRL_{\text{OOC}}\). Более эффективной из двух карт будет та карта, которая дает наименьшую \(MRL_{\text{OOC}}\) для любого заданного сдвига. Из Рисунка 2.10 видно, что контрольная карта, соответствующая пунктирной кривой, работает равномерно лучше.

\begin{figure}[!ht]
	\centering
	\includegraphics[width=0.7\linewidth]{"./images/fig-2.9.png"}
	\caption[]{Кривые медианной длины серии (MRL) для двух контрольных карт в зависимости от величины сдвига}
	\label{fig-2.9}
\end{figure}

Хотя \(ARL\) и \(MRL\) указывают на типичную эффективность карты, а распределение длины серии заметно асимметрично, для лучшего понимания эффективности карты необходимо изучить и оценить все распределение длины серии. С этой целью иногда используются графики, похожие на ящичные диаграммы (см. Radson and Boyd, 2005); примеры показаны на Рисунке \ref{fig-2.10}. Каждая ящичная диаграмма показывает среднее значение (\(ARL\)) распределения в виде черной линии, стандартное отклонение (\(SDRL\)) в виде кружка и медиану (\(MRL\)) в виде серой линии внутри прямоугольника. «Усы» продлены до 5-го и 95-го процентилей (отмечены крестиками) вместо обычных минимума и максимума.

Из \ref{fig-2.10}, например, видно, что карты 1 и 2 имеют схожее распределение длины серии, которое отличается от распределения карт 3 и 4 как по положению, так и по разбросу. Обратите также внимание, что в каждом случае значения \(SDRL\) ближе к значениям \(ARL\).

\begin{figure}[!ht]
	\centering
	\includegraphics[width=0.7\linewidth]{"./images/fig-2.10.png"}
	\caption[]{Диаграмма типа «ящик с усами», обобщающая распределение длины серии контрольных карт}
	\label{fig-2.10}
\end{figure}

Недостатки использования \(ARL\) суммированы ниже. Использование \(MRL\) обеспечит гарантированный контроль над FAR в том смысле, что не более 50\% ложных тревог будут реализованы до \(MRL_{\text{IC}}\). Более того, практически нет недостатков в использовании \(MRL\) для описания типичной эффективности карты. Однако обратите внимание, что при использовании \(MRL\) для характеристики типичной эффективности карты или центра распределения длины серии разброс измеряется межквартильным размахом (IQR), который представляет собой разность между третьим и первым квартилями.

При сравнении эффективности карт на основе \(MRL\) номинальная \(MRL_{\text{IC}}\) фиксируется на приемлемо «высоком» уровне, и предпочтительной является карта с меньшей \(MRL_{\text{OOC}}\). С этой целью обратите внимание, что распределение длины серии карты Шухарта для случая K следует геометрическому распределению с вероятностью успеха \(p = \frac{1}{ARL_{\text{IC}}}\). Для стандартных отраслевых значений номинальной \(ARL_{\text{IC}}\), таких как 370, \(p = \frac{1}{370} = 0,0027\). Следовательно, для распределения с вероятностью успеха \(p\) cdf равна \(P(N \leq t) = 1 - (1 - p)^t\), откуда следует, что \(MRL_{\text{IC}}\) равно \(\frac{\ln(0,5)}{\ln(1 - p)}\). Например, для \(p = 0,0027\) \(MRL_{\text{IC}} = \frac{\ln(0,5)}{\ln(1 - 0,0027)}\), что равно 256,37. В качестве другого примера, для \(p = 0,002\), что соответствует \(ARL_{\text{IC}}\) 500, \(MRL_{\text{IC}}\) равно \(\frac{\ln(0,5)}{\ln(1 - 0,002)} = 346,2\). Поэтому, сохраняя соответствие с традиционной картой Шухарта и для простоты, мы рекомендуем использовать 250 или 350 в качестве целевой номинальной \(MRL_{\text{IC}}\). Другие номинальные значения \(MRL_{\text{IC}}\), соответствующие номинальным значениям \(ARL_{\text{IC}}\), таким как 370, могут быть найдены аналогично. Далее мы приводим некоторые комментарии по поводу \(ARL\) и \(MRL\).

\noindent\textbf{Замечание 2.1. Некоторые комментарии об ARL и MRL}

\begin{enumerate}
	\item Распределение длины серии карты сильно асимметрично вправо, и поэтому выводы, основанные на \(ARL\) о типичной эффективности карты, могут вводить в заблуждение.
	\item Сложность интерпретации. Контрольная карта, имеющая, скажем, \(ARL_{\text{IC}}\) 500, будет иметь около 63\% всех длин серий меньше 500 (это можно увидеть, вычислив 63-й процентиль), но 50\% всех длин серий будут меньше 348, поскольку соответствующая \(MRL_{\text{IC}}\) равна 348.
	\item Могут быть проблемы, когда \(ARL\) (среднее значение распределения) математически не существует; см., например, Chakraborti, Van der Laan, and Van de Wiel (2004) и Graham et al. (2012).
	\item Высокое стандартное отклонение. Опять же, поскольку распределение длины серии сильно асимметрично вправо, во многих случаях наблюдается высокое \(SDRL\), что нежелательно. Более того, с практической точки зрения это означает, что могут быть некоторые чрезвычайно большие значения длины серии, и поэтому некоторые характеристики длины серии не могут быть вычислены за разумное время, то есть использование \(ARL\) в качестве показателя эффективности может быть очень трудоемким. Это особенно проблематично в случае IC.
	\item Отсутствие робастности. Как хорошо известно, \(ARL\), будучи средним, не является робастной мерой, и, следовательно, в некоторых приложениях может потребоваться использование усеченного или винзорированного среднего. Заинтересованный читатель отсылается к литературе для более подробной информации.
\end{enumerate}

В заключение мы выступаем за использование процентилей, таких как медиана, при оценке эффективности карты. Для более всестороннего рассмотрения распределения длины серии следует рассмотреть изучение нескольких показателей эффективности, включая \(ARL\), \(SDRL\) и процентили. Возможно, график, похожий на ящичную диаграмму, как упоминалось ранее, может быть полезен в этом отношении для быстрой визуализации.

\subsubsection{Среднее число выборок до сигнала (ANSS)}

Среднее число выборок до сигнала (ANSS) определяется (см. Reynolds and Stoumbos, 2000) как ожидаемое количество выборок по \(n\) наблюдений, взятых с заданного момента времени до момента, когда карта подаст сигнал. Обычно для ANSS (а также для ANOS и ATS; обсуждаются ниже) измерения берутся с начала процесса, то есть с момента времени ноль. Подобно \(ARL\) и \(MRL\), для эффективной контрольной карты \(ANSS_{\text{IC}}\) должна быть большой, а \(ANSS_{\text{OOC}}\) — маленькой. Чтобы сэкономить место, графики опущены; однако при сравнении двух или более карт \(ANSS_{\text{IC}}\) фиксируется на приемлемо высоком уровне, и контрольная карта с наименьшей \(ANSS_{\text{OOC}}\) для некоторого заданного сдвига (\(\delta\)) параметра процесса работает лучше всего.

\subsubsection{Среднее число наблюдений до сигнала (ANOS)}

Среднее число наблюдений до сигнала (ANOS) определяется (см. Reynolds and Stoumbos, 2000) как ожидаемое количество индивидуальных наблюдений, взятых с заданного момента времени до момента, когда карта подаст сигнал. Существует следующее важное соотношение между ANOS и ANSS, а именно: \(ANOS = (n)(ANSS)\). Подобно ANSS, для эффективной контрольной карты \(ANOS_{\text{IC}}\) должна быть большой, а \(ANOS_{\text{OOC}}\) — маленькой. Опять же, чтобы сэкономить место, графики опущены; однако при сравнении двух или более карт \(ANOS_{\text{IC}}\) фиксируется на приемлемо высоком уровне, и контрольная карта с наименьшей \(ANOS_{\text{OOC}}\) для некоторого заданного сдвига (\(\delta\)) параметра процесса работает лучше всего.

\subsubsection{Среднее время до сигнала (ATS)}

Среднее время до сигнала (ATS) — это количество периодов времени, которые проходят до генерации сигнала на контрольной карте (см. Khoo, 2004). ATS необходимо, поскольку ANSS и ANOS могут быть недостаточными, если существуют некоторые периоды времени, в течение которых мы не проводим инспекцию. Для выборок, которые берутся через равные промежутки времени длительностью \(l\),

\[
ATS = ARL \times l.
\]

Если мы проводим 100\% инспекцию, \(n = l\) и, следовательно, \(ATS = ARL \times n\) (см., например, Wu, Luo, and Zhang, 2006). Важно отметить, что \(l\) определяется как длительность периода, с которым мы имеем дело. Этот период может относиться к дню, неделе, часу и т.д. Эти результаты основаны на фиксированных интервалах выборки. Формально мы можем определить \(ATS\) как ожидаемую продолжительность времени до подачи сигнала картой, измеряемую с начала процесса мониторинга. Так, если мы работали в днях, ATS — это среднее количество дней до первого обнаружения сигнала.

Нам нужно учитывать, находится ли процесс в состоянии IC или OOC, когда мы получаем наш \(ATS\). Для IC-процесса со всеми параметрами, либо известными, либо заданными нам, мы имеем

\[
ATS_{\text{IC}} = ARL_{\text{IC}} \times l
\]

так что

\[
ATS_{\text{IC}} = \frac{1}{FAR} \times l.
\]

Аналогично, для OOC-процесса мы имеем

\[
ATS_{\text{OOC}} = \frac{1}{P(\text{Сигнал за период})} \times l.
\]

De Vries and Conlin (2003) определяют целевой \(ATS\) как желаемую IC-эффективность, которую мы хотим. В качестве примера, целевой \(ATS\) в 365 дней и длительность периода в 7 дней подразумевают целевую частоту ошибок I рода \(\frac{l}{\text{Целевой } ATS} = \frac{7}{365} = 0,0192\). Обратите внимание, что это выше, чем номинальная \(FAR = 0,0027\).

Подобно \(ARL\), для эффективной контрольной карты \(ATS_{\text{IC}}\) должна быть большой, а \(ATS_{\text{OOC}}\) — маленькой. Опять же, чтобы сэкономить место, графики опущены; однако при сравнении двух или более карт \(ATS_{\text{IC}}\) фиксируется на приемлемо высоком уровне, и контрольная карта с наименьшей \(ATS_{\text{OOC}}\) для некоторого заданного сдвига (\(\delta\)) параметра процесса работает лучше всего.

\subsubsection{Число проверенных отдельных единиц (I)}

Это просто еще один способ представления результатов, которые мы получили для \(ARL\). Здесь мы выражаем \(ARL\) через количество отдельных единиц, которые были проверены. Если мы предположим, что у нас есть выборки размером \(n\), мы получим, что количество проверенных единиц — это просто произведение среднего числа взятых выборок и объема выборки. Таким образом, мы имеем \(I = (ARL)(n)\).

\subsection{Операционные характеристики (OC-кривые)}

Операционная характеристика (OC) — это вероятность того, что контрольная карта не обнаружит сдвиг в процессе. Необнаружение сдвига рассматривается как событие, когда карта не подает сигнал при первом последующем наблюдении после сдвига. Операционная характеристика — это вероятностная мера того, насколько быстро изменение в процессе может быть обнаружено контрольной картой.

Теперь мы определяем операционную характеристику как \(P(\text{Отсутствие сигнала} \mid \text{Сдвиг размера } \delta) = \beta\). Операционная характеристика строится в зависимости от \(\delta\), что дает нам кривую операционной характеристики (OC-кривую). OC-кривая полезна для изучения влияния размера сдвига на эффективность карты. Очевидно, что OC-кривые также могут с большим успехом использоваться для сравнения различных карт. Для удовлетворительной карты мы хотели бы видеть, что большему сдвигу соответствует меньшая вероятность того, что статистика для нанесения на карту будет отображаться между \(LCL\) и \(UCL\), то есть не подаст сигнал. И наоборот, вероятность сигнала \((1 - \beta)\) также может быть построена.

На Рисунке \ref{fig-2.11} мы построили две OC-кривые, соответствующие двум разным контрольным картам, в зависимости от сдвигов разного размера. Для обеих карт обратите внимание, что мы начинаем с \(\beta = 1\) для \(\delta = 0\). Это было ожидаемо, поскольку мы определенно хотели, чтобы процесс не подавал сигнал, если сдвиг не произошел (т.е. если процесс все еще IC). По мере увеличения \(\delta\) значения \(\beta\) уменьшаются для обеих карт. Таким образом, мы видим, что по мере увеличения сдвига параметра процесса вероятность неполучения сигнала быстро уменьшается; это желательно. При сравнении двух карт на основе их OC-кривых мы находим, что, поскольку пунктирная кривая всегда лежит ниже сплошной кривой, первая карта равномерно лучше.

\begin{figure}[!ht]
	\centering
	\includegraphics[width=0.7\linewidth]{"./images/fig-2.11.png"}
	\caption[]{Операционные характеристики (OC-кривые) для двух контрольных карт в зависимости от величины сдвига}
	\label{fig-2.11}
\end{figure}

Чтобы сделать тот же вывод другим способом, для тех же двух карт построены и показаны две вероятности сигнала на Рисунке \ref{fig-2.12}. Для обеих карт видно, что вероятность наблюдения сигнала возрастает по мере увеличения размера сдвига параметра процесса, что желательно. Теперь, поскольку пунктирная кривая доминирует над сплошной кривой, что согласуется с тем, что мы сказали для Рисунка \ref{fig-2.11}, соответствующая карта равномерно лучше.

\begin{figure}[!ht]
	\centering
	\includegraphics[width=0.7\linewidth]{"./images/fig-2.12.png"}
	\caption[]{Операционные характеристики (OC-кривые) для двух контрольных карт в зависимости от величины сдвига}
	\label{fig-2.12}
\end{figure}

Теперь, когда различные показатели эффективности обсуждены, мы обсудим два практических случая или ситуации, в которых применяются и изучаются контрольные карты, а именно: когда параметры процесса, представляющие интерес, известны или заданы, пока процесс IC, и когда они неизвестны и должны быть оценены до начала мониторинга процесса.

\subsection{Проектирование контрольных карт}

После того как определен параметр, подлежащий мониторингу, и выбрана статистика для нанесения на карту, нам нужно выбрать объем выборки \(n\), частоту выборки и форму контрольных границ. Это называется \textbf{проектированием} контрольной карты. Теперь мы обсудим различные компоненты проектирования карты.

\subsubsection{Объем выборки, частота выборки и переменные объемы выборки}

Объем выборки и частота выборки (длительность времени между двумя последовательными выборками) должны быть указаны до того, как карта может быть внедрена. Они зависят от таких соображений, как время и стоимость, а также от размера сдвига, который нужно обнаружить. Оба могут влиять на статистическую эффективность карты. Большие выборки могут обнаруживать меньшие сдвиги, но они более дороги и требуют больше времени, тогда как меньшие выборки могут быть достаточными, если ожидаемый сдвиг велик.

Частота выборки, то есть как часто собирать данные, зависит от условий на местах, так сказать, то есть опять же от времени и стоимости, а также ожидаемого сдвига. Очевидно, желательно брать большие выборки как можно чаще, но это, вероятно, неосуществимо. Обычно мы можем брать меньшие выборки через более короткие интервалы времени (чаще) или большие выборки через более длительные интервалы времени (реже). Текущая промышленная практика отдает предпочтение первому подходу. Автоматизированные средства сбора данных сделали этот процесс намного проще, настолько, что нередко проводят выборку и проверку каждого произведенного изделия.

Соответствующий объем выборки может быть определен путем изучения лежащей в основе изменчивости процесса и размера сдвига. Как правило, процесс с относительно большой собственной вариацией требует разумно больших выборок для обнаружения OOC-состояния, в отличие от процесса с относительно малой собственной вариацией. В целом, также легче обнаружить небольшой сдвиг с использованием больших выборок, в отличие от случая, когда сдвиг процесса относительно велик, для которого меньшие выборки были бы достаточными.

Однако частота выборки также важна. Взятие больших выборок очень часто является, возможно, наиболее желательным курсом действий для обнаружения любого изменения процесса. Однако это экономически невыгодно. Следовательно, проблема решения об объеме выборки и частоте выборки становится проблемой распределения усилий по выборке, то есть мы либо берем небольшие выборки чаще (через более короткие интервалы времени), либо берем большие выборки реже (через более длительные интервалы времени).

На практике для карт по количественным признакам часто используются выборки объемом от 1 до 25, причем \(n = 5\) является очень популярным выбором. Несмотря на это, иногда конкретный процесс может накладывать свои собственные ограничения на объем выборки и вынуждать \(n = 1\).

Иногда может случиться так, что некоторые выборки имеют разное количество наблюдений, например, хотя \(n = 5\), может оказаться, что в некоторые моменты времени из-за некоторых практических ограничений выборку пришлось взять с \(n = 4\). В этом случае мы имеем \textbf{переменные объемы выборки}, и обычно есть три варианта.

Можно использовать \textbf{средний объем выборки} для расчета контрольных границ, что приводит к набору приблизительных, но постоянных контрольных границ. При использовании этих контрольных границ предполагается, что будущие объемы выборок не сильно отличаются от предыдущих объемов выборок. В конце концов, эти контрольные границы должны применяться ко всем подгруппам. Этот подход лучше всего работает с большими выборками и когда объемы выборок не отличаются более чем на 25\% от среднего объема выборки. Однако, поскольку контрольные карты по количественным признакам обычно основаны на малых выборках, изменения на 25\% или более, следовательно, не являются маловероятными. Подгруппа размером 3, например, представляет собой 25\%-ное уменьшение объема выборки, если средний объем выборки равен \(n = 4\).

Проблемы возникают, когда точка отображается близко к любой из приблизительных контрольных границ. Точка, которая отображается за пределами одной из приблизительных контрольных границ, вполне может оказаться внутри, если мы рассчитаем точные контрольные границы. Возможно и обратное, то есть точка, которая отображается внутри приблизительных контрольных границ, может фактически оказаться за пределами, если мы рассчитаем точные контрольные границы.

\noindent\textbf{Пример 2.2. Контрольная карта Шухарта \(\overline{X}\) для переменных объемов выборки}

Предположим, что столбец (a) Таблицы \ref{tab:2.2} представляет индивидуальные измерения 15 выборок, которые были смоделированы из нормально распределенного процесса с известным IC-средним \(\mu_0 = 5\) и известным IC-стандартным отклонением \(\sigma_0 = 2\). Однако обратите внимание, что объемы выборок (\(n_i\)) для этих выборок, как показано в столбце (b), различны и должны учитываться при расчете контрольных границ. Если мы собираемся использовать метод среднего объема выборки, нам нужно найти средний объем выборки, то есть \(\overline{n} = \frac{1}{15}\sum_{i=1}^{15} n_i = 4,33 \approx 4\). Следовательно, набор приблизительных (но постоянных) 3-сигма контрольных границ задается как

\[
UCL = \mu_0 + 3\frac{\sigma_0}{\sqrt{\overline{n}}} = 5 + 3\frac{2}{\sqrt{4}} = 8
\]

и

\[
LCL = \mu_0 - 3\frac{\sigma_0}{\sqrt{\overline{n}}} = 5 - 3\frac{2}{\sqrt{4}} = 2
\]

с центральной линией \(CL = \mu_0 = 5\).

\begin{table}[htbp]
	\centering
	\caption{Измерения из процесса с переменными объемами выборки.}
	\label{tab:2.2}
	\renewcommand{\arraystretch}{1.1} % Немного увеличиваем высоту строк для читаемости
	\begin{tabular}{c *{5}{r} c r}
		\toprule
		Выборка & \multicolumn{5}{c}{(a)} & (b) & (c) \\
		\cmidrule(lr){2-6}
		\(i\) & \(X_1\) & \(X_2\) & \(X_3\) & \(X_4\) & \(X_5\) & \(n_i\) & \(\overline{X}_i\) \\
		\midrule
		1  & 0.46 & 2.80 & 7.87 & ---   & ---   & 3 & 3.710 \\
		2  & 9.68 & 5.60 & 2.10 & 9.14  & 2.02  & 5 & 5.708 \\
		3  & 6.18 & 4.93 & 5.27 & 7.54  & ---   & 4 & 5.980 \\
		4  & 3.94 & 5.53 & 5.32 & 4.33  & 4.63  & 5 & 4.750 \\
		5  & 6.43 & 7.93 & 3.92 & 8.64  & ---   & 4 & 6.730 \\
		6  & 5.81 & 0.89 & 5.50 & 2.85  & 3.98  & 5 & 3.806 \\
		7  & 1.96 & 6.00 & 5.65 & ---   & ---   & 3 & 4.537 \\
		8  & 1.37 & 1.28 & 7.17 & 3.74  & 6.93  & 5 & 4.098 \\
		9  & 7.94 & 4.01 & 4.69 & 7.65  & ---   & 4 & 6.073 \\
		10 & 8.33 & 5.75 & 6.01 & 4.57  & ---   & 4 & 6.165 \\
		11 & 3.59 & 1.41 & 2.34 & 9.71  & 4.57  & 5 & 4.324 \\
		12 & 4.11 & 7.43 & 3.50 & 7.68  & 0.66  & 5 & 4.676 \\
		13 & 8.77 & 5.53 & 6.37 & 9.50  & ---   & 4 & 7.543 \\
		14 & 4.97 & 6.00 & 8.62 & 2.73  & ---   & 4 & 5.580 \\
		15 & 4.81 & 3.40 & 8.29 & 5.20  & 5.37  & 5 & 5.414 \\
		\bottomrule
	\end{tabular}
\end{table}

На Рисунке \ref{fig-2.13} показаны выборочные средние \(\overline{X}_i\) из столбца (c), нанесенные на контрольную карту Шухарта \(\overline{X}\) вместе с этими приблизительными контрольными границами.

\begin{figure}[!ht]
	\centering
	\includegraphics[width=0.7\linewidth]{"./images/fig-2.13.png"}
	\caption[]{Операционные характеристики (OC-кривые) для двух контрольных карт в зависимости от величины сдвига}
	\label{fig-2.13}
\end{figure}

На Рисунке \ref{fig-2.13} все точки отображаются между контрольными границами, и в данных нет аномальных паттернов. Следовательно, мы заключаем, что процесс функционирует в состоянии статистической управляемости.

\subsubsection{Переменные контрольные границы}

Хотя предыдущий подход к переменным объемам выборки упрощает расчеты и интерпретацию контрольных карт, с доступностью современных компьютерных технологий естественно попытаться рассчитать точные контрольные границы на основе фактического размера \(n_i\) каждой выборки. Здесь контрольные границы рассчитываются отдельно для каждой подгруппы на основе объема подгруппы \(n_i\). Это означает, что разброс между \(UCL\) и \(LCL\) будет меняться (увеличиваться или уменьшаться) по мере изменения объема выборки. Однако изменяющаяся ширина — это лишь indication различного объема информации, доступной из каждой выборки, как показано различными объемами выборок. Следовательно, визуальное представление этих контрольных границ не так привлекательно, как у границ, основанных на среднем объеме выборки с постоянными контрольными границами. С одной стороны, наличие фиксированных, а не переменных контрольных границ, безусловно, более приятно для глаза. С другой стороны, интерпретация этих переменных контрольных границ остается той же. Если точка отображается на или выше UCL, или на или ниже LCL, это рассматривается как сигнал возможного OOC-процесса, тогда как точка, которая отображается между двумя контрольными границами, рассматривается как отсутствие сигнала и связана с IC-процессом.

\noindent\textbf{Пример 2.2 (продолжение)}

Для контрольной карты Шухарта \(\overline{X}\) с точными контрольными границами, в случае переменных объемов выборки, различные объемы выборок (\(n_i\)), как показано в столбце (b) Таблицы \ref{tab:2.3}, учитываются при расчете контрольных границ. Для этих данных 3-сигма переменные контрольные границы можно найти из \(LCL_i / UCL_i = \mu_0 \pm 3\frac{\sigma_0}{\sqrt{n_i}}\) для \(i = 1, 2, 3, \ldots, 15\), и они представлены в столбцах (d) и (e) Таблицы \ref{tab:2.3}.

Чтобы проиллюстрировать расчеты, рассмотрим выборку номер 1 с \(n_1 = 3\). 3-сигма переменные контрольные границы рассчитываются следующим образом:

\[
LCL_1 = 5 - 3\frac{2}{\sqrt{3}} = 1,536 \quad \text{и} \quad UCL_1 = 5 + 3\frac{2}{\sqrt{3}} = 8,464.
\]

Тогда как для выборки номер 2 с \(n_2 = 5\) 3-сигма переменные контрольные границы рассчитываются следующим образом:

\[
LCL_2 = 5 - 3\frac{2}{\sqrt{5}} = 2,317 \quad \text{и} \quad UCL_2 = 5 + 3\frac{2}{\sqrt{5}} = 7,683.
\]

\begin{table}[htbp]
	\centering
	\caption{Измерения из процесса с переменными объемами выборки.}
	\label{tab:2.3}
	\renewcommand{\arraystretch}{1.1} % Немного увеличиваем высоту строк
	% Используем *{10}{c} для 10 столбцов с выравниванием по центру
	\begin{tabular}{c *{5}{r} c c r r}
		\toprule
		Выборка & \multicolumn{5}{c}{(a)} & (b) & (c) & (d) & (e) \\
		\cmidrule(lr){2-6}
		\(i\) & \(X_1\) & \(X_2\) & \(X_3\) & \(X_4\) & \(X_5\) & \(n_i\) & \(\overline{X}_i\) & \(LCL_i\) & \(UCL_i\) \\
		\midrule
		1  & 0.46 & 2.80 & 7.87 & ---   & ---   & 3 & 3.710 & 1.536 & 8.464 \\
		2  & 9.68 & 5.60 & 2.10 & 9.14  & 2.02  & 5 & 5.708 & 2.317 & 7.683 \\
		3  & 6.18 & 4.93 & 5.27 & 7.54  & ---   & 4 & 5.980 & 2.000 & 8.000 \\
		4  & 3.94 & 5.53 & 5.32 & 4.33  & 4.63  & 5 & 4.750 & 2.317 & 7.683 \\
		5  & 6.43 & 7.93 & 3.92 & 8.64  & ---   & 4 & 6.730 & 2.000 & 8.000 \\
		6  & 5.81 & 0.89 & 5.50 & 2.85  & 3.98  & 5 & 3.806 & 2.317 & 7.683 \\
		7  & 1.96 & 6.00 & 5.65 & ---   & ---   & 3 & 4.537 & 1.536 & 8.464 \\
		8  & 1.37 & 1.28 & 7.17 & 3.74  & 6.93  & 5 & 4.098 & 2.317 & 7.683 \\
		9  & 7.94 & 4.01 & 4.69 & 7.65  & ---   & 4 & 6.073 & 2.000 & 8.000 \\
		10 & 8.33 & 5.75 & 6.01 & 4.57  & ---   & 4 & 6.165 & 2.000 & 8.000 \\
		11 & 3.59 & 1.41 & 2.34 & 9.71  & 4.57  & 5 & 4.324 & 2.317 & 7.683 \\
		12 & 4.11 & 7.43 & 3.50 & 7.68  & 0.66  & 5 & 4.676 & 2.317 & 7.683 \\
		13 & 8.77 & 5.53 & 6.37 & 9.50  & ---   & 4 & 7.543 & 2.000 & 8.000 \\
		14 & 4.97 & 6.00 & 8.62 & 2.73  & ---   & 4 & 5.580 & 2.000 & 8.000 \\
		15 & 4.81 & 3.40 & 8.29 & 5.20  & 5.37  & 5 & 5.414 & 2.317 & 7.683 \\
		\bottomrule
	\end{tabular}
\end{table}

На Рисунке \ref{fig-2.14} показаны выборочные средние \(\overline{X}_i\) из столбца (c), нанесенные на контрольную карту Шухарта \(\overline{X}\) вместе с этими переменными контрольными границами. Как видно из Рисунка \ref{fig-2.14}, расстояние между контрольными границами расширяется и сужается в соответствии с конкретным объемом выборки или подгруппы.

\begin{figure}[!ht]
	\centering
	\includegraphics[width=0.7\linewidth]{"./images/fig-2.14.png"}
	\caption[]{Операционные характеристики (OC-кривые) для двух контрольных карт в зависимости от величины сдвига}
	\label{fig-2.14}
\end{figure}

На Рисунке \ref{fig-2.14} все точки отображаются между контрольными границами, и в данных нет аномальных паттернов. Следовательно, мы заключаем, что процесс функционирует в состоянии статистической управляемости.

\subsubsection{Стандартизированные контрольные границы}

Третий подход к решению проблемы переменных объемов выборки состоит в использовании стандартизированных статистик для нанесения на карту, что приводит к стандартизированным контрольным границам. Пусть \(T_i\) — статистика для нанесения на карту для \(i\)-й подгруппы с IC-средним \(\mu_{T,\text{IC}}\) и IC-стандартным отклонением \(\sigma_{T,\text{IC}}\). Стандартизированная статистика для нанесения на карту для подгруппы \(i\) находится путем вычисления \(Z_i = (T_i - \mu_{T,\text{IC}}) / \sigma_{T,\text{IC}}\), \(i = 1, 2, 3, \ldots\). Предполагается, что эти стандартизированные статистики для нанесения на карту приблизительно следуют стандартному нормальному распределению, то есть \(Z_i \sim N(0,1)\), и, следовательно, \(LCL\) и \(UCL\) можно просто принять равными \(-3\) и \(+3\) соответственно, если мы используем 3-сигма подход для расчета контрольных границ.

Дополнительные усилия по расчету стандартизированных карт компенсируются главным преимуществом, которое теперь у нас есть: постоянные контрольные границы, независимо от объема выборки для подгрупп. Однако возможная трудность с интерпретацией стандартизированных статистик для нанесения на карту или стандартизированных контрольных карт заключается в том, что они больше не находятся в исходной шкале измерений.

\begin{table}[htbp]
	\centering
	\caption{Измерения из процесса с переменными объемами выборки.}
	\label{tab:2.4}
	\renewcommand{\arraystretch}{1.1} % Немного увеличиваем высоту строк
	\begin{tabular}{c *{5}{r} c c r}
		\toprule
		Sample & \multicolumn{5}{c}{(a)} & (b) & (c) & (d) \\
		\cmidrule(lr){2-6}
		\(i\) & \(X_1\) & \(X_2\) & \(X_3\) & \(X_4\) & \(X_5\) & \(n_i\) & \(\overline{X}_i\) & \(Z_i\) \\
		\midrule
		1  & 0.46 & 2.80 & 7.87 & ---   & ---   & 3 & 3.710 & -1.117 \\
		2  & 9.68 & 5.60 & 2.10 & 9.14  & 2.02  & 5 & 5.708 & 0.792  \\
		3  & 6.18 & 4.93 & 5.27 & 7.54  & ---   & 4 & 5.980 & 0.980  \\
		4  & 3.94 & 5.53 & 5.32 & 4.33  & 4.63  & 5 & 4.750 & -0.280 \\
		5  & 6.43 & 7.93 & 3.92 & 8.64  & ---   & 4 & 6.730 & 1.730  \\
		6  & 5.81 & 0.89 & 5.50 & 2.85  & 3.98  & 5 & 3.806 & -1.335 \\
		7  & 1.96 & 6.00 & 5.65 & ---   & ---   & 3 & 4.537 & -0.401 \\
		8  & 1.37 & 1.28 & 7.17 & 3.74  & 6.93  & 5 & 4.098 & -1.008 \\
		9  & 7.94 & 4.01 & 4.69 & 7.65  & ---   & 4 & 6.073 & 1.073  \\
		10 & 8.33 & 5.75 & 6.01 & 4.57  & ---   & 4 & 6.165 & 1.165  \\
		11 & 3.59 & 1.41 & 2.34 & 9.71  & 4.57  & 5 & 4.324 & -0.756 \\
		12 & 4.11 & 7.43 & 3.50 & 7.68  & 0.66  & 5 & 4.676 & -0.362 \\
		13 & 8.77 & 5.53 & 6.37 & 9.50  & ---   & 4 & 7.543 & 2.543  \\
		14 & 4.97 & 6.00 & 8.62 & 2.73  & ---   & 4 & 5.580 & 0.580  \\
		15 & 4.81 & 3.40 & 8.29 & 5.20  & 5.37  & 5 & 5.414 & 0.463  \\
		\bottomrule
	\end{tabular}
\end{table}

\noindent\textbf{Пример 2.2 (продолжение)}

Если мы используем стандартизированные статистики для нанесения на карту из столбца (d) Таблицы \ref{tab:2.4}, мы снова получаем набор постоянных контрольных границ. Чтобы проиллюстрировать расчеты, рассмотрим выборку номер 1 с \(n_1 = 3\). Первая статистика для нанесения на карту, которая является стандартизированным выборочным средним, рассчитывается следующим образом:

\[
Z_1 = \frac{\overline{X}_1 - \mu_0}{\sigma_0 / \sqrt{n_1}} = \frac{3,710 - 5}{2 / \sqrt{3}} = -1,117.
\]

На Рисунке \ref{fig-2.15} показаны стандартизированные выборочные средние \(Z_i\) из столбца (d), нанесенные на контрольную карту Шухарта \(\overline{X}\) вместе с этими стандартизированными контрольными границами.

\begin{figure}[!ht]
	\centering
	\includegraphics[width=0.7\linewidth]{"./images/fig-2.15.png"}
	\caption[]{Операционные характеристики (OC-кривые) для двух контрольных карт в зависимости от величины сдвига}
	\label{fig-2.15}
\end{figure}

На Рисунке \ref{fig-2.15} все точки отображаются между контрольными границами, и в данных нет аномальных паттернов. Следовательно, мы заключаем, что процесс функционирует в состоянии статистической управляемости.

\subsection{Размер сдвига}

Размер сдвига относится к абсолютной величине (или расстоянию), на которое изменился отслеживаемый параметр процесса, при этом принимая во внимание изменчивость процесса. Например, предположим, что мы отслеживаем среднее значение процесса, IC-среднее которого равно \(\mu_0\), а IC-стандартное отклонение равно \(\sigma_0\), и мы заинтересованы в обнаружении того, претерпевает ли среднее процесса устойчивый сдвиг к новому значению \(\mu_1\). Тогда принято выражать «новое» среднее процесса (\(\mu_1\)) через «старое» IC-среднее процесса (\(\mu_0\)) как \(\mu_1 = \mu_0 + \delta\sigma_0\). Таким образом, абсолютный размер сдвига среднего процесса равен \(|\mu_1 - \mu_0| = |\delta\sigma_0|\), тогда как относительный размер сдвига среднего процесса равен \(\delta = |\mu_1 - \mu_0| / \sigma_0\), выраженный в единицах стандартного отклонения.

Вместо работы с абсолютным сдвигом среднего процесса удобнее работать с относительным сдвигом среднего процесса. Например, предположим, что мы отслеживаем два независимых процесса (A и B) с одинаковым IC-средним процесса, \(\mu_A = \mu_B = 10\), скажем, но процесс A и процесс B имеют разные стандартные отклонения, \(\sigma_A = 1\) и \(\sigma_B = 2\), скажем. Если мы пытаемся обнаружить, изменилось ли среднее процесса A до 11, а среднее процесса B до 12, абсолютные размеры сдвигов средних процессов различны, то есть \(|11 - 10| = 1\) для процесса A и \(|12 - 10| = 2\) для процесса B. Однако относительные размеры сдвигов средних процессов одинаковы, то есть \(\delta = |11 - 10| / 1 = |12 - 10| / 2 = 1\).

Таким образом, у нас было два разных процесса (A и B), и мы пытались обнаружить два seemingly разных размера сдвига, но выражение этих сдвигов в единицах стандартного отклонения привело нас к выводу, что мы на самом деле пытаемся обнаружить одинаковый размер сдвига — первый является абсолютным размером сдвига, а второй — относительным размером сдвига. Следовательно, в дальнейшем мы будем называть относительный размер сдвига просто размером сдвига и специально упоминать, когда мы имеем в виду абсолютный размер сдвига.

Когда \(\delta\) меньше 2,0, принято называть сдвиг параметра процесса малым сдвигом, тогда как для больших значений \(\delta\) сдвиг считается большим.

\subsection{Выбор контрольных границ}

Важным соображением при внедрении контрольной карты является выбор констант для нанесения на карту. Для карты Шухарта константой для нанесения на карту является \(k\), и выбор \(k\) напрямую влияет на статистические свойства карты. Увеличение \(k\), например, или перемещение контрольных границ дальше от \(CL\), уменьшает вероятность объявления процесса OOC (когда процесс IC), то есть уменьшает FAR, но в то же время расширение границ увеличивает вероятность не объявить процесс OOC (когда процесс действительно OOC), то есть увеличивает \(\beta\)-риск. И наоборот, перемещение контрольных границ ближе к \(CL\) будет иметь противоположный эффект. То есть в этом случае FAR увеличивается, а \(\beta\)-риск уменьшается.

\subsubsection{\(k\)-сигма границы}

Мы представили \(k\)-сигма границы Шухарта ранее в этой главе. С этими границами мы можем рассчитать FAR (\(\alpha\)) и \(\beta\)-риск (\(\beta\)) для любого выбора \(k\). Однако, поскольку контрольные границы на типичной контрольной карте Шухарта обычно устанавливаются на расстоянии трех стандартных отклонений от \(CL\) (среднего), в качестве примера мы сосредоточим наше внимание на этом конкретном выборе контрольных границ. Напомним, что «3-сигма» контрольные границы задаются как

\[
UCL = \mu_T + 3\sigma_T
\]
\[
CL = \mu_T
\]
\[
LCL = \mu_T - 3\sigma_T.
\]

Предполагая, что выборочное распределение статистики для нанесения на карту \(T\) является нормальным или приблизительно нормальным, мы можем рассчитать FAR (т.е. \(\alpha\)). Однако удобнее рассчитать дополнительную вероятность \(1 - \alpha\), которая является вероятностью отсутствия сигнала, когда процесс IC. Обратите внимание, что

\[
1 - \alpha = P(LCL < T < UCL) = P(\mu_T - 3\sigma_T < T < \mu_T + 3\sigma_T) = P(-3 < Z < 3) = 0,9973.
\]

Таким образом, FAR или вероятность ошибки I рода (\(\alpha\)), которая является вероятностью того, что статистика для нанесения на карту \(T\) выйдет за пределы 3-сигма контрольных границ, равна \(1 - 0,9973 = 0,0027\), когда процесс фактически IC. Это означает, что мы будем наблюдать в среднем неправильный OOC-сигнал или ложную тревогу в 27 из 10 000 или примерно в 1 из каждых 370 выборок из процесса.

\subsubsection{Вероятностные границы}

В качестве альтернативы использованию фиксированной или заданной константы \(k\) (например, 3) в контрольных границах Шухарта, что приводит к некоторой FAR (например, 0,0027) при предполагаемом распределении (нормальном), можно быть заинтересованным в задании FAR и определении соответствующего \(k\) (и, следовательно, контрольных границ). Такие контрольные границы называются \textbf{вероятностными границами} и особенно ценны, когда выборочное распределение статистики для нанесения на карту, используемой для мониторинга интересующего параметра процесса, является асимметричным.

Если мы хотим FAR 0,0027, например, контрольные границы Шухарта должны быть симметрично размещены на расстоянии \(k\)-сигма единиц от \(CL\), где \(k\) находится из

\[
P(T \geq \mu_T + k\sigma_T) = P(Z \geq k) = 0,00135
\]

или эквивалентно из

\[
P(T \leq \mu_T - k\sigma_T) = P(Z \leq -k) = 0,00135
\]

с \(Z \sim N(0,1)\). Таким образом, требуемая константа \(k\) — это просто 99,865-й процентиль IC-распределения стандартизированной статистики \(Z\), что дает нам 99,73\%-ные вероятностные границы. Любое из этих уравнений дает \(k = 3,0\), и результирующая карта является знакомой 3-сигма контрольной картой Шухарта. Следовательно, 99,73\%-ные FAR-вероятностные границы совпадают с 3-сигма границами. Аналогично, для FAR 0,01 мы находим \(k = 2,58\), так что 99\%-ные FAR-вероятностные границы совпадают с 2,58-сигма контрольными границами. Следовательно, вероятностные границы в этом случае однозначно соответствуют определенным \(k\)-сигма границам, и наоборот, и это потому, что IC-распределение \(T\) предполагается нормальным, которое является симметричным. Это не так, однако, когда IC-распределение \(T\) является асимметричным. Например, предположим, что статистика для нанесения на карту \(T\) имеет хи-квадрат распределение с 5 степенями свободы, когда процесс IC. В этом случае можно использовать 95-й и 5-й процентили IC-распределения, чтобы найти 90\%-ные FAR-вероятностные границы; верхние хвостовые вероятности для хи-квадрат распределения приведены в Таблице F Приложения A. Однако можно также использовать 97-й и 7-й процентили хи-квадрат распределения с 5 степенями свободы. Два набора контрольных границ с одинаковой FAR 90\% различны.

Подводя итог, при подходе с фиксированными границами расстояние контрольных границ от \(CL\) задается пользователем, что приводит к определенной FAR (\(\alpha\)), тогда как при вероятностных границах пользователь задает желаемое номинальное значение FAR, которое затем определяет расстояние (\(k\)) контрольных границ от \(CL\).

Можно отметить, что приведенные выше расчеты предполагают подход с равными хвостами (или симметричный), в том смысле, что вероятность того, что статистика для нанесения на карту отобразится выше \(UCL\), равна вероятности того, что она отобразится ниже \(LCL\). Другой способ взглянуть на это — вспомнить тот факт, что \(UCL\) и \(LCL\) оба находятся на расстоянии \(k\) стандартных отклонений от \(CL\). Этот подход кажется разумным, когда распределение статистики для нанесения на карту симметрично. Однако в общем случае можно взять

\[
\alpha_L + \alpha_U = \alpha
\]

где

\[
\alpha_L = P(\text{Статистика для нанесения на карту отображается ниже } LCL \mid \text{Процесс фактически IC})
\]

и

\[
\alpha_U = P(\text{Статистика для нанесения на карту отображается выше } UCL \mid \text{Процесс фактически IC}).
\]

Эти две вероятности называются \(\alpha\)-нижняя (или \(FAR_L\)) и \(\alpha\)-верхняя (или \(FAR_U\)) соответственно. Их задание особенно значимо при использовании статистики для нанесения на карту, которая имеет асимметричное распределение.

В случаях, когда представляет интерес только восходящий или нисходящий сдвиг, так что ищется односторонняя контрольная карта, или когда восходящий (или нисходящий) сдвиг представляет меньший (или больший) интерес, чем нисходящий (восходящий) сдвиг, может быть полезно неравное распределение FAR. В первом случае односторонней карты либо \(\alpha_U = 0\), либо \(\alpha_L = 0\), а в последнем случае \(\alpha_L + \alpha_U = \alpha\), где, например, \(0 < \alpha_U < \alpha_L < 1\). При этом подходе с неравными хвостами (или асимметричном) и предполагая нормальное распределение для статистики для нанесения на карту, нам нужно найти решения уравнений \(P(Z \leq -z_{\alpha_L}) = \alpha_L\) и \(P(Z \geq z_{\alpha_U}) = \alpha_U\), чтобы найти \(LCL\) и \(UCL\) соответственно. Требуемые константы \(-z_{\alpha_L}\) и \(z_{\alpha_U}\) могут быть легко найдены как \(100\alpha_L\)-й и \(100\alpha_U\)-й процентили \(Z\)-распределения соответственно. Другой альтернативой в подобной ситуации может быть рассмотрение односторонней (верхней или нижней) контрольной карты только с верхней (или нижней) контрольной границей, где можно использовать либо фиксированный подход, либо подход на основе вероятностных границ для определения требуемых констант.

\noindent\textbf{Пример 2.3. Неравные хвосты (или асимметричные) контрольные границы}

Предположим, что интересующая статистика для нанесения на карту имеет нормальное распределение, и предположим, что \(\alpha_L\) задана равной 0,0015, тогда как \(\alpha_U\) задана как 0,0005. Используя \(\alpha_L + \alpha_U = \alpha\), мы находим FAR равной 0,002. Решая для \(z_{0,0015}\) и \(z_{0,0005}\) из \(P(Z \leq -z_{0,0015}) = 0,0015\) и из \(P(Z \geq z_{0,0005}) = 0,0005\), мы получаем \(z_{\alpha_L} = z_{0,0015} = 2,97\) и \(z_{\alpha_U} = z_{0,0005} = 3,29\) соответственно. Таким образом, контрольные границы могут быть определены из Уравнения 2.1 как

\[
UCL = \mu_T + 3,29\sigma_T
\]

и

\[
LCL = \mu_T - 2,97\sigma_T.
\]

В этом случае \(LCL\) ближе к \(CL\), чем \(UCL\). Этот подход может использоваться, когда восходящий сдвиг конкретного параметра процесса вызывает большее беспокойство, чем нисходящий сдвиг. Далее рассмотрим одностороннюю верхнюю контрольную карту.

\noindent\textbf{Пример 2.4. Односторонняя верхняя контрольная граница}

Предположим, что нас интересует только обнаружение восходящего сдвига среднего процесса, который следует нормальному распределению. Тогда, если номинальная FAR задана как 0,0025, пусть \(\alpha_L = 0\) и \(\alpha_U = 0,0025\). Затем нам нужно найти \(z_{0,0025}\) из \(P(Z \geq z_{0,0025}) = 0,0025\), что дает \(z_{\alpha_U} = z_{0,0025} = 2,81\), 99,75-й процентиль \(Z\)-распределения. Таким образом, \(UCL\) устанавливается на уровне \(UCL = \mu_T + 2,81\sigma_T\) без \(LCL\) на контрольной карте.

Во многих ситуациях истинное лежащее в основе распределение характеристики качества неизвестно, и, следовательно, выборочное распределение статистики для нанесения на карту также неизвестно. Поэтому расчет точных вероятностных границ проблематичен. Обратите внимание, что в таком случае можно использовать \(k\)-сигма границы для некоторого заданного значения \(k\), но при нахождении FAR этой контрольной карты это может привести к той же проблеме. Однако, если статистика для нанесения на карту является средним или может быть выражена как среднее, приблизительно, можно использовать центральную предельную теорему (которая подразумевает, что среднее \(n\) i.i.d. случайных величин приблизительно следует нормальному распределению независимо от лежащего в основе распределения отдельных переменных). Например, если \(X_1, \ldots, X_n\) являются i.i.d. со средним \(\mu\) и дисперсией $\overline{X} = \frac{1}{n}\sum_{i=1}^n X_i \sim N\left(\mu, \frac{\sigma^2}{n}\right)$ приблизительно, для достаточно больших значений $n$, независимо от распределения генеральной совокупности. Таким образом, например, можно использовать 3-сигма контрольные границы для среднего, которые, как предполагается, дают FAR приблизительно 0,0027. Однако, если истинное лежащее в основе распределение процесса заметно отличается от нормального, свойства контрольной карты (такие как FAR), основанные на приблизительных границах, могут сильно пострадать, особенно поскольку объем выборки \(n\), как правило, невелик в типичном приложении контроля процесса. В таких ситуациях у пользователя есть другие варианты. Некоторые из них обсуждаются далее в этой книге.